\documentclass[12pt,a4paper]{article}

%% ---------- mathematics ----------
\usepackage{amsmath,amsthm,amssymb,mathtools}

%% ---------- graphics & diagrams ----------
\usepackage{tikz}
\usetikzlibrary{arrows.meta,positioning,shapes.geometric,calc,fit,backgrounds,matrix,decorations.pathreplacing}
\usepackage{graphicx}

%% ---------- tables ----------
\usepackage{booktabs,tabularx,multirow,longtable,array}

%% ---------- algorithms ----------
\usepackage[ruled,vlined,linesnumbered]{algorithm2e}

%% ---------- coloured boxes ----------
\usepackage[most]{tcolorbox}

%% ---------- hyperlinks & cross-refs ----------
\usepackage{hyperref}
\usepackage[capitalize,noabbrev]{cleveref}

%% ---------- page layout ----------
\usepackage[margin=1in]{geometry}
\usepackage{fancyhdr}
\usepackage{enumitem}
\usepackage{xcolor}

%% ---------- fonts & misc ----------
\usepackage{microtype}

%% ============================================================
%%  Custom tcolorbox environments
%% ============================================================
\newtcolorbox{keyinsight}[1][]{%
  colback=blue!5!white,colframe=blue!75!black,
  fonttitle=\bfseries,title=Key Insight,#1}

\newtcolorbox{example}[1][]{%
  colback=green!5!white,colframe=green!75!black,
  fonttitle=\bfseries,title=Example,#1}

\newtcolorbox{warningbox}[1][]{%
  colback=red!5!white,colframe=red!75!black,
  fonttitle=\bfseries,title=Warning,#1}

%% ============================================================
%%  amsthm environments
%% ============================================================
\newtheorem{theorem}{Theorem}[section]
\newtheorem{lemma}[theorem]{Lemma}
\newtheorem{definition}{Definition}[section]
\newtheorem{remark}{Remark}[section]
\newtheorem{proposition}[theorem]{Proposition}
\newtheorem{corollary}[theorem]{Corollary}

%% ============================================================
%%  Custom commands (collected from all sessions in this file)
%% ============================================================
\newcommand{\xt}{x_t}
\newcommand{\ht}{h_t}
\newcommand{\hprev}{h_{t-1}}
\newcommand{\sigmoid}{\sigma}
\newcommand{\R}{\mathbb{R}}
\newcommand{\odotop}{\odot}

\pagestyle{fancy}
\fancyhf{}
\lhead{\textsc{Deep Learning} --- IIIT Dharwad}
\rhead{LSTM: Theory, Gates, and Numerical Practice}
\rfoot{Page \thepage}
\lfoot{Prof.\ S R M Prasanna \& Dr.\ Sunil Saumya}
\renewcommand{\headrulewidth}{0.4pt}

\title{%
  {\Large\bfseries Deep Learning Course}\\[6pt]
  {\LARGE LSTM: Theory, Gates, and Numerical Practice}\\[4pt]
  {\normalsize IIIT Dharwad --- M.Tech Programme}
}
\author{Prof.\ S R M Prasanna \& Dr.\ Sunil Saumya}
\date{Sessions 18, 19, 20}

\begin{document}

\maketitle
\tableofcontents
\clearpage

%% ============================================================
%% Session 18: RNN Hands-On and LSTM Introduction
%% ============================================================
\part{Session 18: RNN Hands-On and LSTM Introduction}

%%%%%%%%%%%%%%%%%%%%%%%%%%%%%%%%%%%%%%%%%%%%%%%%%%%%%%%%%%%%%%%%%%%%%%%%%%%%%

%% ── Title page ──────────────────────────────────────────────────────────────
\begin{titlepage}
  \centering
  \vspace*{2cm}
  {\Huge\bfseries Deep Learning Course\\[0.4em]
   Session 18 Lecture Notes\par}
  \vspace{1.2cm}
  {\LARGE\itshape RNN Hands-On Implementation \&\\[0.3em]
   Introduction to LSTM\par}
  \vspace{1.5cm}
  \rule{\linewidth}{0.6pt}\\[0.4cm]
  \begin{tabular}{ll}
    \textbf{Session Number:}  & 18 \\[4pt]
    \textbf{Topic:}           & Recurrent Neural Networks -- hands-on; \\
                              & unrolled RNN as feedforward network; \\
                              & motivation for LSTM \\[4pt]
    \textbf{Date:}            & January 17, 2026 \\[4pt]
    \textbf{Instructors:}     & Prof.\ S R M Prasanna \& Dr.\ Sunil Saumya \\[4pt]
    \textbf{Institution:}     & IIIT Dharwad \\[4pt]
    \textbf{Slide Deck:}      & \texttt{08-PyTorchDL-RNN} (slide 20/29) \\[4pt]
    \textbf{Notebook:}        & \texttt{RNNClassification.ipynb} (Google Colab) \\[4pt]
    \textbf{Dataset:}         & IMDB Movie Reviews (50\,000 samples) \\
                              & Shakespeare Dataset (500 plays) via Kaggle \\
  \end{tabular}
  \rule{\linewidth}{0.6pt}
  \vfill
  {\large These notes synthesise the lecture transcript and visual slide
   extractions for Session~18 of the deep learning course at IIIT Dharwad.
   They cover the architecture of simple RNNs, text pre-processing,
   the PyTorch implementation of a sentiment-analysis RNN, and the
   theoretical motivation for LSTM.\par}
\end{titlepage}

\newpage

%%%%%%%%%%%%%%%%%%%%%%%%%%%%%%%%%%%%%%%%%%%%%%%%%%%%%%%%%%%%%%%%%%%%%%%%%%%%%
\section{Introduction}
\label{sec:introduction}
%%%%%%%%%%%%%%%%%%%%%%%%%%%%%%%%%%%%%%%%%%%%%%%%%%%%%%%%%%%%%%%%%%%%%%%%%%%%%

Session~18 is a pivotal bridge between RNN theory and real-world
implementation.  The session opened with a brief review of L2
regularisation and a Q\&A on assignments, then devoted the majority of
its $\approx 1$h\,52\,min runtime to two interlocking themes:

\begin{enumerate}[label=\textbf{\arabic*.}]
  \item \textbf{RNN as an unrolled feedforward network.}  The instructor
        revisited the computational graph of an RNN unrolled to time step
        $t=3$, establishing the shared weight matrices $W$ and $U$, and
        deriving the hidden-state update equation.
  \item \textbf{Hands-on sentiment analysis.}  A complete PyTorch
        implementation of a simple RNN for IMDB movie-review sentiment
        classification was built from scratch inside a Google Colab
        notebook (\texttt{RNNClassification.ipynb}).  Every stage of the
        pipeline -- data loading, text pre-processing, vocabulary
        construction, dataset/dataloader classes, model definition, and
        training loop -- was coded and explained live.
\end{enumerate}

The session concluded by returning to the unrolled RNN diagram to
explain \emph{why} long sequences create a vanishing-gradient problem,
directly motivating the LSTM architecture to be introduced in the
next session.

\begin{keyinsight}
The central architectural difference between a standard ANN and an RNN
is the \emph{recurrent connection}: when processing word $x_t$, the
network also receives the hidden state $h_{t-1}$ produced when
processing $x_{t-1}$.  This allows the network to maintain a
``memory'' of previous inputs, which is essential for sequential data.
\end{keyinsight}

%%%%%%%%%%%%%%%%%%%%%%%%%%%%%%%%%%%%%%%%%%%%%%%%%%%%%%%%%%%%%%%%%%%%%%%%%%%%%
\section{Background and Prerequisites}
\label{sec:background}
%%%%%%%%%%%%%%%%%%%%%%%%%%%%%%%%%%%%%%%%%%%%%%%%%%%%%%%%%%%%%%%%%%%%%%%%%%%%%

\subsection{Sequential Data and the Limits of ANNs}
\label{subsec:ann_limits}

A standard feed-forward neural network (ANN) treats every input
independently.  Given a vocabulary of size $V$, each word $w$ is
represented as a one-hot vector $\mathbf{x} \in \{0,1\}^V$, and the
network applies:

\begin{equation}
  \mathbf{o} = f\!\left(W_{\text{in}}\,\mathbf{x} + \mathbf{b}\right),
  \label{eq:ann_forward}
\end{equation}

where $f$ is a non-linear activation function.  Two limitations arise
immediately:

\begin{enumerate}[label=(\alph*)]
  \item \textbf{Vocabulary-sized input.}  For $V = 75\,509$ (the vocabulary
        of the IMDB corpus used in this session), $W_{\text{in}}$ has
        $75\,509 \times h$ parameters for a hidden layer of width $h$.
        At $h=64$ that is already $\approx 4.8$\,M weights \emph{just}
        for the first layer.
  \item \textbf{No temporal dependency.}  Processing the words
        ``\emph{not}'' and ``\emph{good}'' simultaneously in a window
        discards the positional relationship between them.  The network
        sees both as part of a fixed-length input vector; it does not
        ``know'' which came first.
\end{enumerate}

Window-based approaches (feeding $k$ consecutive words as a single
input) partially address (b) but cannot handle variable-length
sequences without padding, and they still treat the $k$ words as a
single unordered bag.

\subsection{L2 Regularisation Review}
\label{subsec:l2_review}

The session opened with a clarification about L2 regularisation,
which had appeared in an assignment question.  The standard form adds
a penalty proportional to the squared weight magnitudes to the loss:

\begin{equation}
  \mathcal{L}_{\text{reg}} = \mathcal{L} + \frac{\lambda}{2}
  \sum_{j} w_j^2,
  \label{eq:l2_regularisation}
\end{equation}

where $\lambda$ is the regularisation strength (\texttt{weight\_decay}
in PyTorch optimisers).  The factor of $\tfrac{1}{2}$ is a
mathematical convenience: when differentiating, the $2$ in the
exponent cancels with $\tfrac{1}{2}$, yielding the clean gradient
$\nabla_{w_j}\mathcal{L}_{\text{reg}} = \nabla_{w_j}\mathcal{L} +
\lambda w_j$.

\begin{example}[title=L2 Penalty Calculation]
  Given $\lambda = 10^{-4}$ and a weight $w = 2.0$:
  \begin{align}
    \text{With factor } \tfrac{1}{2}:&\quad
      \frac{\lambda}{2} w^2 = \frac{10^{-4}}{2} \times 4 = 2 \times 10^{-4} \\
    \text{Without factor } \tfrac{1}{2}:&\quad
      \lambda w^2 = 10^{-4} \times 4 = 4 \times 10^{-4}
  \end{align}
  Both forms are mathematically valid; the choice must be consistent
  with the convention used when defining the loss function.
\end{example}

%%%%%%%%%%%%%%%%%%%%%%%%%%%%%%%%%%%%%%%%%%%%%%%%%%%%%%%%%%%%%%%%%%%%%%%%%%%%%
\section{RNN Architecture: Unrolled View}
\label{sec:rnn_architecture}
%%%%%%%%%%%%%%%%%%%%%%%%%%%%%%%%%%%%%%%%%%%%%%%%%%%%%%%%%%%%%%%%%%%%%%%%%%%%%

\subsection{Hidden-State Update Equation}
\label{subsec:hidden_state}

At every time step $t$, a simple RNN computes:

\begin{equation}
  \boxed{h_t = f\!\left(W\,h_{t-1} + U\,x_t + b\right)}
  \label{eq:rnn_update}
\end{equation}

where:
\begin{itemize}[noitemsep]
  \item $x_t \in \mathbb{R}^{d}$ is the embedding of the word at
        position $t$ ($d = 100$ in the notebook);
  \item $h_t \in \mathbb{R}^{n_h}$ is the hidden state ($n_h = 64$
        neurons);
  \item $W \in \mathbb{R}^{n_h \times n_h}$ is the
        \emph{hidden-to-hidden} weight matrix (shared across all $t$);
  \item $U \in \mathbb{R}^{n_h \times d}$ is the
        \emph{input-to-hidden} weight matrix (shared across all $t$);
  \item $b \in \mathbb{R}^{n_h}$ is the bias vector;
  \item $f$ is a point-wise non-linearity (typically $\tanh$).
\end{itemize}

The output at time step $t=3$ as shown on slide V01 is:

\begin{equation}
  O_3 = f\!\left(W \cdot O_2 + U \cdot X_3\right),
  \label{eq:o3}
\end{equation}

which is the equation explicitly displayed on the course slide.

\begin{keyinsight}[title=Parameter Sharing is the Key]
  The matrices $W$ and $U$ are \textbf{identical} at every time step.
  This is why an RNN can handle sequences of any length without
  increasing the number of parameters.  A 391-word review and a
  61-word review both use the same $W$ and $U$; only the number of
  times they are applied differs.
\end{keyinsight}

\subsection{Unrolled RNN as a Deep Feedforward Network}
\label{subsec:unrolled}

\Cref{fig:rnn_unrolled} shows the unrolled computational graph for a
sequence of length $T$.  Each column corresponds to one time step;
the rows represent layers.  The horizontal arrows represent the flow
of hidden state from $h_{t-1}$ to $h_t$, while the vertical arrows
represent the local computation within each time step.

\begin{figure}[ht]
  \centering
  \begin{tikzpicture}[
      node distance=1.8cm and 2.2cm,
      box/.style={draw,rounded corners=3pt,minimum width=1.4cm,
                  minimum height=0.8cm,fill=blue!10,font=\small},
      inp/.style={draw,circle,minimum size=0.7cm,fill=orange!20,font=\small},
      out/.style={draw,circle,minimum size=0.7cm,fill=green!20,font=\small},
      arr/.style={-{Stealth[length=5pt]},thick},
      rarr/.style={-{Stealth[length=5pt]},thick,blue!70!black}
    ]

    %% Hidden state nodes
    \node[box] (h0) {$h_0$};
    \node[box,right=of h0] (h1) {$h_1$};
    \node[box,right=of h1] (h2) {$h_2$};
    \node[box,right=of h2] (h3) {$h_3$};

    %% Recurrent arrows
    \draw[rarr] (h0) -- node[above,font=\tiny]{$W$} (h1);
    \draw[rarr] (h1) -- node[above,font=\tiny]{$W$} (h2);
    \draw[rarr] (h2) -- node[above,font=\tiny]{$W$} (h3);

    %% Input nodes (below)
    \node[inp,below=1.2cm of h0] (x0) {$x_0$};
    \node[inp,below=1.2cm of h1] (x1) {$x_1$};
    \node[inp,below=1.2cm of h2] (x2) {$x_2$};
    \node[inp,below=1.2cm of h3] (x3) {$x_3$};

    %% Input arrows
    \draw[arr] (x0) -- node[right,font=\tiny]{$U$} (h0);
    \draw[arr] (x1) -- node[right,font=\tiny]{$U$} (h1);
    \draw[arr] (x2) -- node[right,font=\tiny]{$U$} (h2);
    \draw[arr] (x3) -- node[right,font=\tiny]{$U$} (h3);

    %% Output nodes (above)
    \node[out,above=1.2cm of h0] (o0) {$O_0$};
    \node[out,above=1.2cm of h1] (o1) {$O_1$};
    \node[out,above=1.2cm of h2] (o2) {$O_2$};
    \node[out,above=1.2cm of h3] (o3) {$O_3$};

    %% Output arrows
    \draw[arr] (h0) -- (o0);
    \draw[arr] (h1) -- (o1);
    \draw[arr] (h2) -- (o2);
    \draw[arr] (h3) -- (o3);

    %% Labels
    \node[font=\footnotesize,below=0.2cm of x0] {$t=0$};
    \node[font=\footnotesize,below=0.2cm of x1] {$t=1$};
    \node[font=\footnotesize,below=0.2cm of x2] {$t=2$};
    \node[font=\footnotesize,below=0.2cm of x3] {$t=3$};

    %% Initial state h_{-1}
    \node[left=1.2cm of h0,font=\small] (hinit) {$h_{-1}=\mathbf{0}$};
    \draw[rarr] (hinit) -- node[above,font=\tiny]{$W$} (h0);

  \end{tikzpicture}
  \caption{RNN unrolled over four time steps $t = 0, 1, 2, 3$.  The
           weight matrices $W$ (hidden-to-hidden) and $U$
           (input-to-hidden) are \emph{shared} at every column.
           The highlighted output $O_3$ at $t=3$ satisfies
           $O_3 = f(W \cdot O_2 + U \cdot X_3)$, as shown on slide V01.}
  \label{fig:rnn_unrolled}
\end{figure}

\begin{remark}
  The unrolled view makes clear that an RNN with sequence length $T$
  is equivalent to a feedforward network of depth $T$.  For a 391-word
  IMDB review, the unrolled network has 391 ``layers'', all sharing
  the same parameters $W$, $U$.  This depth is precisely why
  gradients can vanish during backpropagation through time (BPTT).
\end{remark}

\subsection{ANN vs.\ RNN: Architectural Comparison}
\label{subsec:ann_vs_rnn}

\Cref{tab:ann_vs_rnn} summarises the key differences between a
standard ANN and an RNN for sequential NLP tasks, as discussed by
the instructor during the session.

\begin{table}[ht]
  \centering
  \caption{ANN versus RNN for sequential text processing.}
  \label{tab:ann_vs_rnn}
  \begin{tabularx}{\linewidth}{lXX}
    \toprule
    \textbf{Property} & \textbf{ANN} & \textbf{RNN} \\
    \midrule
    Input at time $t$ & Fixed-size window of $k$ words
                      & Single word $x_t$; also receives $h_{t-1}$ \\
    \addlinespace
    Sequence order & Implicit (position in window)
                   & Explicit (strictly sequential) \\
    \addlinespace
    Variable-length input & Requires padding to fixed $k$
                          & Handled naturally (no padding needed) \\
    \addlinespace
    Parameter count & $V \times h$ for first layer ($V = 75\,509$)
                    & $d \times n_h + n_h \times n_h$ (after embedding) \\
    \addlinespace
    Long-range context & Limited to window size $k$
                       & Theoretically unlimited (in practice limited by vanishing gradient) \\
    \addlinespace
    Parallelism & All words in window processed simultaneously
                & Sequential -- $t+1$ requires $h_t$ \\
    \addlinespace
    Typical use & Classification with fixed context
                & Sequence classification, language modelling \\
    \bottomrule
  \end{tabularx}
\end{table}

%%%%%%%%%%%%%%%%%%%%%%%%%%%%%%%%%%%%%%%%%%%%%%%%%%%%%%%%%%%%%%%%%%%%%%%%%%%%%
\section{Text Pre-Processing Pipeline}
\label{sec:preprocessing}
%%%%%%%%%%%%%%%%%%%%%%%%%%%%%%%%%%%%%%%%%%%%%%%%%%%%%%%%%%%%%%%%%%%%%%%%%%%%%

The IMDB dataset contains 50\,000 labelled movie reviews (25\,000
positive, 25\,000 negative).  The session used a randomly sampled
subset of 5\,000 reviews to keep training times manageable on Colab.
\Cref{fig:preprocessing_pipeline} illustrates the five-stage pipeline.

\begin{figure}[ht]
  \centering
  \begin{tikzpicture}[
      stage/.style={draw,rectangle,rounded corners=4pt,
                    minimum width=2.4cm,minimum height=0.9cm,
                    fill=teal!15,font=\small,align=center,thick},
      arr/.style={-{Stealth[length=6pt]},thick,teal!60!black}
    ]
    \node[stage] (s1) {Raw\\Text};
    \node[stage,right=1.3cm of s1] (s2) {Lower-\\case};
    \node[stage,right=1.3cm of s2] (s3) {Remove\\Punct.};
    \node[stage,right=1.3cm of s3] (s4) {Tokenise};
    \node[stage,right=1.3cm of s4] (s5) {Build\\Vocab};
    \node[stage,below=1.2cm of s5] (s6) {Text\\$\to$ Indices};
    \node[stage,below=1.2cm of s3] (s7) {Tensor\\Dataset};
    \node[stage,below=1.2cm of s1] (s8) {Data\\Loader};

    \draw[arr] (s1) -- (s2);
    \draw[arr] (s2) -- (s3);
    \draw[arr] (s3) -- (s4);
    \draw[arr] (s4) -- (s5);
    \draw[arr] (s5) -- (s6);
    \draw[arr] (s6) -- (s7);
    \draw[arr] (s7) -- (s8);

    %% labels below arrows
    \node[font=\tiny,below=0.1cm of s2] {step 1};
    \node[font=\tiny,below=0.1cm of s3] {step 2};
    \node[font=\tiny,below=0.1cm of s4] {step 3};
    \node[font=\tiny,below=0.1cm of s5] {step 4};
  \end{tikzpicture}
  \caption{Five-stage text pre-processing pipeline implemented in
           \texttt{RNNClassification.ipynb}.}
  \label{fig:preprocessing_pipeline}
\end{figure}

\subsection{Tokenisation and Cleaning}
\label{subsec:tokenisation}

The \texttt{tokenize(text)} function applied the following
transformations in sequence:

\begin{enumerate}[label=\arabic*.]
  \item Convert to lower-case.
  \item Remove all punctuation characters using the \texttt{string.punctuation} set.
  \item Split on whitespace to produce a list of tokens.
\end{enumerate}

These three steps are illustrated below in Python pseudocode:

\begin{verbatim}
import string

def tokenize(text):
    text = text.lower()
    text = ''.join(c for c in text if c not in string.punctuation)
    tokens = text.split()
    return tokens
\end{verbatim}

\begin{example}[title=Tokenisation in Action]
  Input: \texttt{``One of the most beautifully crafted films you'll ever see.''}\\[4pt]
  After lower-case: \texttt{``one of the most beautifully crafted films youll ever see''}\\[4pt]
  After removing punctuation: same (apostrophe removed; ``you'll'' $\to$ ``youll'')\\[4pt]
  After splitting: \texttt{[`one', `of', `the', `most', `beautifully', `crafted', `films', `youll', `ever', `see']}
\end{example}

Note that the apostrophe in ``you'll'' is stripped, making the
token \texttt{youll} (index~190 in the vocabulary).  A student asked
why this was not classified as the \texttt{<UNK>} token -- the answer
is that tokenisation removes the apostrophe \emph{before} vocabulary
lookup, so \texttt{youll} appears as a known token in the vocabulary.

\subsection{Vocabulary Construction}
\label{subsec:vocabulary}

The vocabulary is a Python dictionary mapping each unique token to an
integer index:

\begin{equation}
  \text{vocab} = \bigl\{\,\text{token} \mapsto \text{index}\,\bigr\}
  \label{eq:vocab}
\end{equation}

Two special tokens are pre-inserted before any text is processed:

\begin{center}
  \begin{tabular}{cll}
    \toprule
    \textbf{Index} & \textbf{Token} & \textbf{Purpose} \\
    \midrule
    0 & \texttt{<PAD>} & Padding (needed when batches contain variable-length sequences) \\
    1 & \texttt{<UNK>} & Unknown token (used at inference for out-of-vocabulary words) \\
    \bottomrule
  \end{tabular}
\end{center}

Every new unique token encountered while iterating through the 5\,000
training reviews is assigned the next available index
($|\text{vocab}|$ at the time of insertion):

\begin{verbatim}
def build_vocab(texts):
    vocab = {'<PAD>': 0, '<UNK>': 1}
    for text in texts:
        for token in tokenize(text):
            if token not in vocab:
                vocab[token] = len(vocab)
    return vocab
\end{verbatim}

After processing all 5\,000 rows the vocabulary contained
$|\text{vocab}| = 75\,509$ unique tokens.  The instructor noted that
processing the full 50\,000-sample dataset would likely exceed
100\,000 unique tokens.

\begin{keyinsight}[title=Why Both \texttt{<PAD>} and \texttt{<UNK>}?]
  \textbf{\texttt{<UNK>}} (index 1) handles out-of-vocabulary words at
  test time.  If a review contains a word not seen during training, it
  is mapped to index~1.  \textbf{\texttt{<PAD>}} (index 0) is included
  as a standard convention even though this notebook processes each
  review individually (batch size = 1), so no padding is actually
  applied.  Padding becomes necessary when batching multiple
  variable-length sequences together.
\end{keyinsight}

\subsection{Text-to-Indices Conversion}
\label{subsec:text_to_indices}

The function \texttt{text\_to\_indices(text, vocab)} converts a raw
review string into a tensor of vocabulary indices:

\begin{verbatim}
def text_to_indices(text, vocab):
    tokens = tokenize(text)
    return [vocab[t] if t in vocab else vocab['<UNK>'] for t in tokens]
\end{verbatim}

The sequence of indices preserves the word order of the original
review.  For example, the sentence ``one of the \ldots'' maps to the
integer sequence $[162, 12, 16, \ldots]$.  This is the input
$[x_0, x_1, x_2, \ldots]$ that will be fed token-by-token to the RNN.

\subsection{Padding and Batch Processing Trade-offs}
\label{subsec:padding}

A student raised the question of padding during the live session.
\Cref{tab:padding_tradeoffs} summarises the trade-offs discussed:

\begin{table}[ht]
  \centering
  \caption{Trade-offs between padded batching and individual sequence processing.}
  \label{tab:padding_tradeoffs}
  \begin{tabularx}{\linewidth}{lXX}
    \toprule
    \textbf{Property} & \textbf{Batch with Padding} & \textbf{Batch Size = 1} \\
    \midrule
    GPU utilisation & Higher (parallelism) & Lower \\
    \addlinespace
    Information loss & Yes -- zeros dilute gradient signal & None \\
    \addlinespace
    Memory & Fixed (max-length sequences) & Variable \\
    \addlinespace
    Implementation & Requires \texttt{<PAD>} masking & Simple \\
    \addlinespace
    Recommended for RNN & Only when resource-constrained & Preferred \\
    \bottomrule
  \end{tabularx}
\end{table}

\begin{warningbox}
  Padding with zeros is never ideal for RNNs.  The padded time steps
  produce meaningless hidden-state updates, dilute the gradient signal
  during backpropagation, and can degrade model performance.  The
  best practice is to process each sequence individually
  (\texttt{batch\_size = 1}) or to use \texttt{pack\_padded\_sequence}
  in PyTorch, which skips padded positions entirely.
\end{warningbox}

%%%%%%%%%%%%%%%%%%%%%%%%%%%%%%%%%%%%%%%%%%%%%%%%%%%%%%%%%%%%%%%%%%%%%%%%%%%%%
\section{PyTorch Dataset and DataLoader}
\label{sec:dataset_dataloader}
%%%%%%%%%%%%%%%%%%%%%%%%%%%%%%%%%%%%%%%%%%%%%%%%%%%%%%%%%%%%%%%%%%%%%%%%%%%%%

\subsection{Custom Dataset Class}
\label{subsec:custom_dataset}

Following the standard PyTorch pattern introduced in earlier sessions,
a custom \texttt{IMDBDataset} class was defined by inheriting from
\texttt{torch.utils.data.Dataset} and implementing three methods:

\begin{itemize}[noitemsep]
  \item \texttt{\_\_init\_\_(self, df, vocab)}: Stores the DataFrame
        (with a reset index) and the vocabulary dictionary.
  \item \texttt{\_\_len\_\_(self)}: Returns the number of reviews.
  \item \texttt{\_\_getitem\_\_(self, idx)}: Converts the $idx$-th
        review to a tensor of indices and returns the
        (review\_tensor, label\_tensor) pair.
\end{itemize}

The \texttt{\_\_getitem\_\_} method calls \texttt{text\_to\_indices}
on-the-fly, so no pre-encoded copy of the entire dataset needs to be
held in memory.  The label is returned as a floating-point tensor
($0.0$ for negative, $1.0$ for positive) to match the binary
cross-entropy loss expectation.

\begin{verbatim}
class IMDBDataset(Dataset):
    def __init__(self, df, vocab):
        self.df = df.reset_index(drop=True)
        self.vocab = vocab

    def __len__(self):
        return len(self.df)

    def __getitem__(self, idx):
        review = self.df.loc[idx, 'review']
        label  = self.df.loc[idx, 'sentiment']
        x = torch.tensor(text_to_indices(review, self.vocab))
        y = torch.tensor(float(label))
        return x, y
\end{verbatim}

The reason for \texttt{reset\_index(drop=True)} is that the 5\,000
rows were drawn by \emph{random sampling} from the original 50\,000,
so the index values are non-contiguous (e.g.\ 4, 312, 7003, \ldots).
Resetting creates a clean 0-to-4999 integer index that \texttt{iloc}
can use directly.

\subsection{DataLoader Configuration}
\label{subsec:dataloader}

The DataLoader was configured with \texttt{batch\_size = 1} for the
reasons discussed in \cref{subsec:padding}:

\begin{verbatim}
train_loader = DataLoader(train_dataset,
                          batch_size=1,
                          shuffle=True)
\end{verbatim}

Iterating the loader yields batches of shape
$1 \times T_i$ for the review tensor (where $T_i$ is the number of
tokens in the $i$-th review) and shape $1$ for the label scalar.  The
instructor demonstrated this live, showing:

\begin{itemize}[noitemsep]
  \item Batch 0: review shape $1 \times 61$ (61 tokens) -- label: positive.
  \item Batch 1: review shape $1 \times 391$ (391 tokens) -- label: negative.
  \item Batch 2: review shape $1 \times 173$ (173 tokens).
\end{itemize}

This confirms that RNNs handle variable-length inputs naturally -- the
DataLoader does not need to pad or truncate any sequence.

%%%%%%%%%%%%%%%%%%%%%%%%%%%%%%%%%%%%%%%%%%%%%%%%%%%%%%%%%%%%%%%%%%%%%%%%%%%%%
\section{Simple RNN Model Implementation}
\label{sec:model}
%%%%%%%%%%%%%%%%%%%%%%%%%%%%%%%%%%%%%%%%%%%%%%%%%%%%%%%%%%%%%%%%%%%%%%%%%%%%%

\subsection{Model Architecture Overview}
\label{subsec:model_overview}

The \texttt{SimpleRNN} model consists of exactly three learnable
components stacked in sequence:

\begin{enumerate}[label=\arabic*.]
  \item \textbf{Embedding layer} (\texttt{nn.Embedding}): maps each
        integer token index to a dense $d$-dimensional vector.
  \item \textbf{RNN layer} (\texttt{nn.RNN}): processes the embedded
        sequence token-by-token, maintaining a hidden state.
  \item \textbf{Linear (fully-connected) layer} (\texttt{nn.Linear}):
        maps the final hidden state to a scalar sentiment score.
\end{enumerate}

\Cref{fig:model_architecture} illustrates the full forward pass.

\begin{figure}[ht]
  \centering
  \begin{tikzpicture}[
      layer/.style={draw,rectangle,rounded corners=5pt,
                    minimum width=3.5cm,minimum height=0.9cm,
                    fill=violet!12,font=\small,align=center,thick},
      data/.style={draw,rectangle,rounded corners=3pt,
                   minimum width=3.5cm,minimum height=0.7cm,
                   fill=gray!12,font=\small,align=center},
      arr/.style={-{Stealth[length=6pt]},thick}
    ]

    \node[data]  (inp)   {Token indices $[i_0,i_1,\ldots,i_{T-1}]$};
    \node[layer,below=0.8cm of inp]  (emb)
        {Embedding\\$V \times 100 \to T \times 100$};
    \node[layer,below=0.8cm of emb]  (rnn)
        {RNN\\input=100, hidden=64\\(batch\_first=True)};
    \node[data,below=0.6cm of rnn]   (hn)
        {$h_T \in \mathbb{R}^{64}$ (final hidden state)};
    \node[layer,below=0.6cm of hn]   (fc)
        {Linear $64 \to 1$};
    \node[data,below=0.6cm of fc]    (out)
        {Scalar logit $\hat{y} \in \mathbb{R}$};
    \node[layer,below=0.6cm of out,fill=orange!15]  (loss)
        {BCEWithLogitsLoss\\$\mathcal{L}(\hat{y}, y)$};

    \draw[arr] (inp)  -- (emb);
    \draw[arr] (emb)  -- node[right,font=\tiny]{embedded sequence} (rnn);
    \draw[arr] (rnn)  -- node[right,font=\tiny]{output, $h_n$} (hn);
    \draw[arr] (hn)   -- (fc);
    \draw[arr] (fc)   -- (out);
    \draw[arr] (out)  -- (loss);

  \end{tikzpicture}
  \caption{Forward pass of the \texttt{SimpleRNN} model used in the
           session.  $V = 75\,509$ (vocabulary size), $T$ = sequence
           length (variable), $d = 100$ (embedding dimension),
           $n_h = 64$ (RNN hidden size).}
  \label{fig:model_architecture}
\end{figure}

\subsection{The Embedding Layer}
\label{subsec:embedding}

The embedding layer performs two roles simultaneously:

\begin{enumerate}[label=(\alph*)]
  \item \textbf{One-hot encoding:} internally, the layer treats each
        integer index as selecting the corresponding row from a learnable
        weight matrix $E \in \mathbb{R}^{V \times d}$.  This is
        mathematically identical to multiplying a one-hot vector by $E$.
  \item \textbf{Dimensionality reduction:} the 75\,509-dimensional
        one-hot vector is compressed to a dense $d=100$-dimensional
        vector.
\end{enumerate}

\begin{equation}
  \mathbf{e}_t = E[i_t] \in \mathbb{R}^{d}, \quad d = 100
  \label{eq:embedding}
\end{equation}

The two key advantages of using embeddings rather than raw one-hot
vectors:

\begin{enumerate}[noitemsep]
  \item \textbf{Parameter reduction:} the input-to-hidden weight
        matrix changes from $V \times n_h = 75\,509 \times 64
        \approx 4.8$\,M parameters to $d \times n_h = 100 \times 64
        = 6\,400$ parameters -- a reduction by a factor of $\approx 750$.
  \item \textbf{Denser gradients:} a one-hot vector with 75\,508 zeros
        produces almost no gradient signal.  A dense embedding vector
        with non-zero values everywhere propagates gradients more
        effectively.
\end{enumerate}

\begin{keyinsight}[title=Embedding Layer as Fully-Connected Layer]
  Conceptually, the embedding layer is simply a fully-connected layer
  (without bias or activation) placed \emph{before} the main
  architecture starts.  Its weights $E$ are initialised randomly and
  learned by gradient descent along with all other parameters.  If
  pre-trained word embeddings (Word2Vec, GloVe, FastText) are loaded
  into $E$, the model can leverage semantic information from large
  external corpora.
\end{keyinsight}

\subsection{The RNN Layer}
\label{subsec:rnn_layer}

The \texttt{nn.RNN} module implements \cref{eq:rnn_update} for every
time step in a vectorised manner.  Key constructor arguments used:

\begin{center}
  \begin{tabular}{lll}
    \toprule
    \textbf{Argument} & \textbf{Value} & \textbf{Meaning} \\
    \midrule
    \texttt{input\_size}  & 100  & Dimension of each embedded token \\
    \texttt{hidden\_size} & 64   & Width of the hidden state $h_t$ \\
    \texttt{batch\_first} & \texttt{True} & Input tensor shape: $(B, T, d)$ not $(T, B, d)$ \\
    \bottomrule
  \end{tabular}
\end{center}

When called, \texttt{nn.RNN} returns two tensors:
\begin{itemize}[noitemsep]
  \item \texttt{output}: shape $(B, T, n_h)$ -- the hidden state at
        \emph{every} time step.
  \item \texttt{h\_n}: shape $(1, B, n_h)$ -- the hidden state at the
        \emph{last} time step only.
\end{itemize}

For sentiment classification, only \texttt{h\_n} is needed because
the final hidden state encodes a summary of the entire sequence.
The instructor verified this by showing that the last row of
\texttt{output} is numerically identical to \texttt{h\_n}.

\subsection{Parameter Count}
\label{subsec:param_count}

The total number of trainable parameters in the model is:

\begin{align}
  \text{Embedding:} &\quad V \times d = 75\,509 \times 100 = 7\,550\,900 \label{eq:emb_params} \\
  \text{RNN} \;(U): &\quad d \times n_h = 100 \times 64 = 6\,400 \label{eq:rnn_u} \\
  \text{RNN} \;(W): &\quad n_h \times n_h = 64 \times 64 = 4\,096 \label{eq:rnn_w} \\
  \text{RNN biases:}&\quad 2 \times n_h = 128 \label{eq:rnn_bias} \\
  \text{Linear:}    &\quad n_h \times 1 + 1 = 65 \label{eq:fc_params} \\
  \hline \nonumber \\
  \text{Total:}     &\quad \approx 7\,561\,589 \label{eq:total_params}
\end{align}

The embedding layer dominates the parameter count.  This is typical
for NLP models with large vocabularies; pre-training embeddings or
using sub-word tokenisation (BPE) are common mitigations.

\subsection{Forward Pass Code}
\label{subsec:forward_pass}

\begin{verbatim}
class SimpleRNN(nn.Module):
    def __init__(self, vocab_size, embed_dim, hidden_size, output_size):
        super(SimpleRNN, self).__init__()
        self.embedding = nn.Embedding(vocab_size, embed_dim)
        self.rnn = nn.RNN(embed_dim, hidden_size, batch_first=True)
        self.fc  = nn.Linear(hidden_size, output_size)

    def forward(self, x):
        embedded = self.embedding(x)          # (B, T, 100)
        _, h_n   = self.rnn(embedded)         # h_n: (1, B, 64)
        h_n      = h_n.squeeze(0)             # (B, 64)
        out      = self.fc(h_n)               # (B, 1)
        return out
\end{verbatim}

The underscore \texttt{\_} discards the per-timestep output tensor;
only the final hidden state \texttt{h\_n} is passed to the
fully-connected layer.

%%%%%%%%%%%%%%%%%%%%%%%%%%%%%%%%%%%%%%%%%%%%%%%%%%%%%%%%%%%%%%%%%%%%%%%%%%%%%
\section{Training Procedure}
\label{sec:training}
%%%%%%%%%%%%%%%%%%%%%%%%%%%%%%%%%%%%%%%%%%%%%%%%%%%%%%%%%%%%%%%%%%%%%%%%%%%%%

\subsection{Loss Function and Optimiser}
\label{subsec:loss_optimiser}

Binary cross-entropy with logits was used as the loss function:

\begin{equation}
  \mathcal{L} = -\frac{1}{N}\sum_{i=1}^{N}
    \Bigl[ y_i \log \sigma(\hat{y}_i)
         + (1 - y_i) \log (1 - \sigma(\hat{y}_i)) \Bigr],
  \label{eq:bce_loss}
\end{equation}

where $\sigma$ is the sigmoid function.  \texttt{BCEWithLogitsLoss}
applies the sigmoid internally, which is numerically more stable than
applying it explicitly and then computing the log.

A student asked why a ``linear'' output layer is used instead of
logistic regression.  The answer: the linear layer produces a logit
$\hat{y}$; the \emph{sigmoid} non-linearity is applied inside
\texttt{BCEWithLogitsLoss}, so the full model does implement logistic
regression -- the non-linearity just appears in the loss function rather
than in \texttt{forward()}.

\subsection{Training Loop}
\label{subsec:training_loop}

\Cref{alg:training_loop} presents the training algorithm used in the
session.

\begin{algorithm}[ht]
  \caption{RNN Sentiment Classifier Training Loop}
  \label{alg:training_loop}
  \SetAlgoLined
  \KwIn{Training DataLoader \texttt{train\_loader}, number of epochs
        \texttt{num\_epochs}, model, loss function \texttt{criterion},
        optimiser \texttt{optimizer}}
  \KwOut{Trained model parameters}
  \BlankLine
  model.train()\;
  \For{epoch $\leftarrow 1$ \KwTo \texttt{num\_epochs}}{
    epoch\_loss $\leftarrow 0$\;
    \For{(review, label) \textbf{in} train\_loader}{
      review, label $\leftarrow$ review.to(device), label.to(device)\;
      optimizer.zero\_grad()\;
      output $\leftarrow$ model(review)\tcp*{forward pass}
      output $\leftarrow$ output.squeeze(1)\tcp*{match label shape}
      loss $\leftarrow$ criterion(output, label)\tcp*{compute BCE loss}
      loss.backward()\tcp*{backprop through time (BPTT)}
      optimizer.step()\tcp*{SGD / Adam weight update}
      epoch\_loss $\leftarrow$ epoch\_loss $+$ loss.item()\;
    }
    \textbf{print} epoch, epoch\_loss / len(train\_loader)\;
  }
\end{algorithm}

\subsection{Key Training Observations}
\label{subsec:training_observations}

\begin{itemize}
  \item The model was trained for 10 epochs on the 5\,000-sample subset.
  \item The instructor showed the loss decreasing epoch by epoch,
        confirming that the pipeline was correctly connected.
  \item Because \texttt{batch\_size = 1}, each gradient step is computed
        from a single review -- this is stochastic gradient descent
        (SGD) in its purest form, with very noisy gradient estimates.
  \item Moving to GPU (\texttt{device = 'cuda'} if available) was
        highlighted as essential for larger datasets or more epochs.
\end{itemize}

%%%%%%%%%%%%%%%%%%%%%%%%%%%%%%%%%%%%%%%%%%%%%%%%%%%%%%%%%%%%%%%%%%%%%%%%%%%%%
\section{Hands-On: Shakespeare Next-Character Prediction}
\label{sec:shakespeare}
%%%%%%%%%%%%%%%%%%%%%%%%%%%%%%%%%%%%%%%%%%%%%%%%%%%%%%%%%%%%%%%%%%%%%%%%%%%%%

The visual extractions (V02--V04) also captured code from a
secondary notebook focused on \emph{next-character prediction} using
the Kaggle Shakespeare dataset (500 plays).  This task trains an RNN
to predict the next character $c_{t+1}$ given a sequence of
characters $[c_1, c_2, \ldots, c_t]$.

\subsection{Dataset Download}
\label{subsec:shakespeare_data}

\begin{verbatim}
import kagglehub
path = kagglehub.dataset_download(
    "kylehobbsdev/shakespeare-dataset-of-500-shakespeare-plays"
)
print("Path to dataset files:", path)
# Path to dataset files:
#   /root/.cache/kaggle/datasets/kylehobbsdev/
#   shakespeare-dataset-of-500-shakespeare-plays/versions/1
\end{verbatim}

\subsection{Sliding-Window Sequence Construction}
\label{subsec:sliding_window}

Given the encoded Shakespeare text $[c_1, c_2, \ldots, c_N]$ (where
each $c_i$ is a character index), training pairs are constructed as:

\begin{equation}
  (x^{(k)},\, y^{(k)}) = \bigl([c_k, c_{k+1}, \ldots, c_{k+L-1}],\;
                                 c_{k+L}\bigr), \quad
  k = 1, 2, \ldots, N - L
  \label{eq:sliding_window}
\end{equation}

where $L$ is the context length.  The session notebook constructed
942 such training sequences from the sampled Shakespeare corpus.

\begin{example}[title=Sliding Window Sequences]
  For an encoded text starting with character indices $[3, 7, 5, 2, \ldots]$
  with context length $L = 2$:
  \begin{align*}
    (x^{(1)}, y^{(1)}) &= ([3, 7],\; 5) \\
    (x^{(2)}, y^{(2)}) &= ([3, 7, 5],\; 2) \\
    (x^{(3)}, y^{(3)}) &= ([3, 7, 5, 2],\; \ldots)
  \end{align*}
\end{example}

%%%%%%%%%%%%%%%%%%%%%%%%%%%%%%%%%%%%%%%%%%%%%%%%%%%%%%%%%%%%%%%%%%%%%%%%%%%%%
\section{Vanishing Gradient and LSTM Motivation}
\label{sec:lstm_motivation}
%%%%%%%%%%%%%%%%%%%%%%%%%%%%%%%%%%%%%%%%%%%%%%%%%%%%%%%%%%%%%%%%%%%%%%%%%%%%%

\subsection{Backpropagation Through Time (BPTT)}
\label{subsec:bptt}

Training an RNN requires computing gradients of the loss
$\mathcal{L}$ with respect to the shared weights $W$ and $U$.
Because the hidden state at time $t$ depends on all previous hidden
states, the chain rule produces products of the Jacobian matrices of
$f$:

\begin{equation}
  \frac{\partial \mathcal{L}}{\partial W}
    = \sum_{t=1}^{T} \frac{\partial \mathcal{L}}{\partial h_T}
      \underbrace{\prod_{k=t}^{T-1}
        \frac{\partial h_{k+1}}{\partial h_k}}_{\text{product of }T-t\text{ Jacobians}}
      \frac{\partial h_t}{\partial W}.
  \label{eq:bptt}
\end{equation}

If the spectral norm of each Jacobian $\tfrac{\partial h_{k+1}}{\partial h_k}$
is less than 1 (which commonly occurs with $\tanh$ activations when
$\|W\| < 1$), then for large $T$:

\begin{equation}
  \prod_{k=t}^{T-1} \left\| \frac{\partial h_{k+1}}{\partial h_k} \right\|
  \;\longrightarrow\; 0 \quad \text{as } T - t \to \infty.
  \label{eq:vanishing}
\end{equation}

This is the \textbf{vanishing gradient problem}: gradients for early
time steps become negligibly small, making it impossible to learn
long-range dependencies.

\begin{keyinsight}[title=Unrolled View Explains Vanishing Gradients]
  The unrolled RNN (slide V05) has \emph{depth} equal to sequence
  length $T$.  A 391-word review produces an unrolled network 391
  layers deep.  Just as very deep feedforward networks suffer from
  vanishing gradients in their early layers, very long RNN sequences
  lose gradient signal for early time steps.  LSTM (Long Short-Term
  Memory) solves this by introducing a \emph{cell state} and
  \emph{gating mechanisms} that allow gradients to flow through long
  sequences without vanishing.
\end{keyinsight}

\subsection{Properties of Long Sequences in RNNs}
\label{subsec:long_seq}

Two further issues arise with long sequences:

\begin{enumerate}[label=\arabic*.]
  \item \textbf{Slow processing.}  Because computation is strictly
        sequential (time step $t+1$ waits for $t$), a 500-word review
        requires 500 sequential matrix-vector multiplications.
        Compare this to a Transformer, which processes all positions
        in parallel.
  \item \textbf{Context degradation.}  Even without vanishing
        gradients, the hidden state $h_t$ is a fixed-size vector
        ($n_h = 64$) that must encode the entire history
        $[x_0, \ldots, x_{t-1}]$.  For very long sequences, the
        ``memory'' of the earliest tokens is overwritten by later ones.
\end{enumerate}

The instructor addressed a student's question about building a
question-answering system from a 100-page PDF using RNNs:

\begin{itemize}
  \item Possible in principle: convert the PDF to a structured dataset
        (e.g.\ one row per paragraph), then train a sequence-to-sequence
        RNN.
  \item Not ideal: RNNs struggle with the very long context of an entire
        book, and modern systems (RAG + large language models) perform
        significantly better on such tasks.
\end{itemize}

%%%%%%%%%%%%%%%%%%%%%%%%%%%%%%%%%%%%%%%%%%%%%%%%%%%%%%%%%%%%%%%%%%%%%%%%%%%%%
\section{Comparison of RNN Variants}
\label{sec:rnn_variants}
%%%%%%%%%%%%%%%%%%%%%%%%%%%%%%%%%%%%%%%%%%%%%%%%%%%%%%%%%%%%%%%%%%%%%%%%%%%%%

\Cref{tab:rnn_variants} compares the vanilla RNN introduced in this
session with the more advanced architectures (LSTM and GRU) mentioned
as upcoming topics.

\begin{table}[ht]
  \centering
  \caption{Comparison of RNN, LSTM, and GRU architectures.}
  \label{tab:rnn_variants}
  \begin{tabularx}{\linewidth}{lXXX}
    \toprule
    \textbf{Property} & \textbf{Vanilla RNN} & \textbf{LSTM} & \textbf{GRU} \\
    \midrule
    Hidden state & $h_t$ only & $h_t$ + cell state $c_t$ & $h_t$ only \\
    \addlinespace
    Gates & None & Input, forget, output (3) & Reset, update (2) \\
    \addlinespace
    Vanishing gradient & Severe for long $T$ & Mitigated by cell state & Partially mitigated \\
    \addlinespace
    Parameters (per layer) & $d n_h + n_h^2 + 2n_h$ & $4(d n_h + n_h^2 + 2n_h)$ & $3(d n_h + n_h^2 + 2n_h)$ \\
    \addlinespace
    Training speed & Fastest & Slowest & Intermediate \\
    \addlinespace
    Long-range dependency & Poor & Good & Good \\
    \addlinespace
    Best for & Short sequences & Long sequences & Long sequences \\
    \bottomrule
  \end{tabularx}
\end{table}

\begin{remark}
  For the IMDB sentiment task, both LSTM and GRU typically achieve
  better accuracy than vanilla RNN because movie reviews can be
  hundreds of words long and sentiment often depends on long-range
  context (e.g.\ a negation at the start of the review).
\end{remark}

%%%%%%%%%%%%%%%%%%%%%%%%%%%%%%%%%%%%%%%%%%%%%%%%%%%%%%%%%%%%%%%%%%%%%%%%%%%%%
\section{Summary and Conclusion}
\label{sec:summary}
%%%%%%%%%%%%%%%%%%%%%%%%%%%%%%%%%%%%%%%%%%%%%%%%%%%%%%%%%%%%%%%%%%%%%%%%%%%%%

\subsection{What Was Covered}
\label{subsec:covered}

Session~18 covered the following material:

\begin{enumerate}[label=\textbf{\arabic*.}]
  \item \textbf{L2 Regularisation Review.}  Both forms
        ($\lambda \sum w^2$ and $\tfrac{\lambda}{2} \sum w^2$) are
        mathematically valid; the $\tfrac{1}{2}$ factor is a
        convenience that simplifies the gradient.

  \item \textbf{RNN Architectural Motivation.}  The key change from
        ANN to RNN is the recurrent connection that feeds the previous
        hidden state $h_{t-1}$ as an additional input at each time
        step $t$.  This gives the network an explicit ``memory'' of the
        sequence so far.

  \item \textbf{Text Pre-Processing Pipeline.}  A five-stage pipeline
        was implemented: lower-casing, punctuation removal,
        tokenisation, vocabulary construction (with \texttt{<PAD>} and
        \texttt{<UNK>}), and text-to-index mapping.  The IMDB
        5\,000-sample vocabulary contained 75\,509 unique tokens.

  \item \textbf{PyTorch Dataset and DataLoader.}  A custom
        \texttt{IMDBDataset} class was created following the standard
        PyTorch pattern, with \texttt{batch\_size = 1} to avoid padding.

  \item \textbf{SimpleRNN Model.}  The model comprises three layers:
        \texttt{nn.Embedding}($V$, 100),
        \texttt{nn.RNN}(100, 64, \texttt{batch\_first=True}), and
        \texttt{nn.Linear}(64, 1).  The embedding layer converts sparse
        one-hot vectors to dense 100-dimensional representations,
        reducing parameter count by $\approx 750\times$.

  \item \textbf{Forward Pass Mechanics.}  The RNN processes the embedded
        sequence token by token.  After the final token, the hidden
        state $h_T \in \mathbb{R}^{64}$ encodes the entire sequence
        and is fed to the linear layer for binary classification.

  \item \textbf{Training Loop.}  \texttt{BCEWithLogitsLoss} +
        Adam optimiser; 10 epochs on 5\,000 samples.  Loss decreased
        monotonically, confirming correct implementation.

  \item \textbf{Vanishing Gradient and LSTM Motivation.}  The unrolled
        RNN at depth $T$ suffers from vanishing gradients because
        gradient products shrink exponentially.  LSTM addresses this
        with a gated cell state; GRU uses a simplified two-gate design.
\end{enumerate}

\subsection{Key Equations}
\label{subsec:key_equations}

For quick reference, the three core equations from this session:

\begin{align}
  \text{Hidden state update:} &\quad
    h_t = \tanh(W h_{t-1} + U x_t + b)
    \label{eq:hidden_summary} \\
  \text{Output at } t=3: &\quad
    O_3 = f(W \cdot O_2 + U \cdot X_3)
    \label{eq:o3_summary} \\
  \text{Binary cross-entropy loss:} &\quad
    \mathcal{L} = -\bigl[y \log \sigma(\hat{y})
                       + (1-y)\log(1-\sigma(\hat{y}))\bigr]
    \label{eq:bce_summary}
\end{align}

\subsection{Looking Ahead}
\label{subsec:ahead}

The next session will introduce \textbf{Long Short-Term Memory (LSTM)}
networks.  Key topics to anticipate:

\begin{itemize}[noitemsep]
  \item The cell state $c_t$ as a ``conveyor belt'' for long-range information.
  \item The forget gate, input gate, and output gate.
  \item How LSTM mitigates the vanishing gradient problem.
  \item PyTorch implementation using \texttt{nn.LSTM}.
\end{itemize}

Students were encouraged to practise by:
\begin{enumerate}[noitemsep,label=(\roman*)]
  \item Printing the output of every layer in the notebook to verify
        shapes at each stage.
  \item Extending the training to the full 50\,000-sample dataset and
        comparing accuracy.
  \item Replacing the embedding with pre-trained GloVe vectors and
        observing any accuracy improvement.
\end{enumerate}

\begin{keyinsight}[title=The Big Picture]
  The RNN family -- from the vanilla network in this session through
  LSTM and GRU -- all share the fundamental design principle of
  \cref{eq:rnn_update}: a hidden state that is updated at each time
  step by combining the current input with the previous hidden state.
  Every subsequent innovation (LSTM gates, attention mechanisms,
  Transformers) can be understood as a solution to specific limitations
  of this basic recurrence.
\end{keyinsight}

%%%%%%%%%%%%%%%%%%%%%%%%%%%%%%%%%%%%%%%%%%%%%%%%%%%%%%%%%%%%%%%%%%%%%%%%%%%%%
%% References / Resources
%%%%%%%%%%%%%%%%%%%%%%%%%%%%%%%%%%%%%%%%%%%%%%%%%%%%%%%%%%%%%%%%%%%%%%%%%%%%%
\newpage
\section*{Session Resources}
\addcontentsline{toc}{section}{Session Resources}

\begin{itemize}
  \item \textbf{Video lecture:}
        \url{https://vimeo.com/1156130746}
  \item \textbf{Slide deck (PDF):}
        \url{https://learning.iiitdwd.ac.in/mod/pdfprotect/view.php?id=1556}
  \item \textbf{Notebook:} \texttt{RNNClassification.ipynb}
        (shared via course portal; PyTorch implementation of the
        IMDB sentiment classifier)
  \item \textbf{Dataset -- IMDB:}
        \texttt{datasets.load\_dataset("imdb")} via HuggingFace, or
        download from Kaggle
  \item \textbf{Dataset -- Shakespeare:}
        \texttt{kylehobbsdev/shakespeare-dataset-of-500-shakespeare-plays}
        on Kaggle Hub
\end{itemize}

\clearpage

%% ============================================================
%% Session 19: LSTM: Short vs Long-Term Dependency and Gates
%% ============================================================
\part{Session 19: LSTM: Short vs Long-Term Dependency and Gates}

%% ── Title page ───────────────────────────────────────────────────────────────
\begin{titlepage}
  \centering
  \vspace*{2cm}
  {\LARGE\bfseries Deep Learning Course\\[0.4em]
   Indian Institute of Information Technology Dharwad\\[1.5em]}
  \rule{\linewidth}{1pt}\\[0.8em]
  {\huge\bfseries Session 19\\[0.5em]
   Long Short-Term Memory (LSTM):\\[0.3em]
   Short vs Long-Term Dependency and Gates}\\[0.8em]
  \rule{\linewidth}{1pt}\\[1.5em]
  \begin{tabular}{ll}
    \textbf{Session Number:} & 19 \\[0.4em]
    \textbf{Topic:}          & LSTM motivation, vanishing gradient, gate mechanism \\[0.4em]
    \textbf{Date:}           & January 24, 2026 \\[0.4em]
    \textbf{Duration:}       & $\approx 1$h\,39m \\[0.4em]
    \textbf{Instructors:}    & Prof.\ S.\ R.\ Mahadeva Prasanna \\
                             & Dr.\ Sunil Saumya \\[0.4em]
    \textbf{Slide Decks:}    & \texttt{05-DL-RNN} (slides 30, 35, 37/63) \\
                             & \texttt{06-RNNTypesBackpropagation} (slide 20/22) \\
  \end{tabular}
  \vfill
  {\large\itshape Acknowledgement: Colah's Blog --
   ``Understanding LSTM Networks''}\\[1em]
  {\large IIITDwd -- B.Tech Deep Learning, 2025--26}
\end{titlepage}

%% ── Table of contents ────────────────────────────────────────────────────────

\newpage

% =============================================================================
\section{Introduction}
\label{sec:introduction}
% =============================================================================

Session~19 of the IIITDwd Deep Learning course marks the transition from
standard Recurrent Neural Networks (RNNs) to the Long Short-Term Memory
(LSTM) architecture. The session was delivered jointly by
Prof.\ S.\ R.\ Mahadeva Prasanna and Dr.\ Sunil Saumya and draws primarily
from two slide decks: \texttt{05-DL-RNN} and
\texttt{06-RNNTypesBackpropagation}.

The central thesis of the session is that standard RNNs, while capable of
modelling \emph{short-term} sequential dependencies, fundamentally fail on
\emph{long-term} dependencies because gradients vanish as they propagate
backwards through many time steps. The LSTM architecture, proposed by
Hochreiter and Schmidhuber (1997), addresses this failure through two
complementary innovations:

\begin{enumerate}[label=\textbf{\arabic*.}]
  \item A dedicated \textbf{cell-state highway} $C_t$ that carries long-term
        context with only minor, additive interactions --- preventing the
        multiplicative gradient decay that cripples plain RNNs.
  \item A set of three \textbf{gates} (forget, input, output), each
        implemented as a sigmoid neural network, that selectively read,
        write, and erase information on the cell-state line.
\end{enumerate}

\begin{keyinsight}[title=Session Objective]
By the end of Session~19 students should be able to: (i) articulate why
plain RNNs fail on long-range context using the gradient-vanishing argument;
(ii) describe the two-line (cell-state + hidden-state) feedback structure of
LSTM; and (iii) explain the role of each of the three gates at a conceptual
level. Numerical gate computations are deferred to Session~20.
\end{keyinsight}

% =============================================================================
\section{Background and Prerequisites}
\label{sec:background}
% =============================================================================

\subsection{Sequence Modelling and the RNN Recap}
\label{subsec:rnn-recap}

Before LSTM can be appreciated, it is essential to recall the operation of a
vanilla RNN cell. At each time step $t$ the hidden state is updated as

\begin{equation}
  \label{eq:rnn-update}
  h_t = \tanh\!\bigl(W_{hh}\,h_{t-1} + W_{xh}\,x_t + b_h\bigr),
\end{equation}

where $x_t \in \mathbb{R}^{d}$ is the input vector, $h_{t-1} \in
\mathbb{R}^{n}$ is the previous hidden state, $W_{hh} \in \mathbb{R}^{n
\times n}$ and $W_{xh} \in \mathbb{R}^{n \times d}$ are learnable weight
matrices, and $b_h$ is a bias vector. The output at each time step (e.g.\
for language modelling) is

\begin{equation}
  \label{eq:rnn-output}
  \hat{y}_t = \text{softmax}(W_{hy}\,h_t + b_y).
\end{equation}

The feedback loop through $h_{t-1}$ gives the RNN its \emph{memory}: at
time step $t$ the hidden state encodes a compressed representation of all
tokens seen so far. As the instructors emphasised,

\begin{quote}
``The goal [in language modelling] is, given one word or a set of words in a
given sequence, predict the next word.'' \hfill (Dr.\ Sunil Saumya, 00:15)
\end{quote}

\subsection{What Is Short-Term Dependency?}
\label{subsec:short-term}

\begin{definition}[Short-Term Dependency]
  \label{def:short-term}
  A prediction task exhibits \emph{short-term dependency} when the
  relevant context lies within a small, fixed window immediately preceding
  the prediction point. A standard RNN with a shallow hidden state suffices
  for such tasks.
\end{definition}

\begin{example}[title=``The clouds are in the \underline{\hspace{1.5em}}'']
  Given the phrase ``the clouds are in the'', any fluent speaker immediately
  predicts ``sky''. The information required (clouds, in, the) is entirely
  local --- the gap between evidence and prediction is negligible. This is a
  classic short-term dependency, and a vanilla RNN handles it well.
\end{example}

\subsection{What Is Long-Term Dependency?}
\label{subsec:long-term}

\begin{definition}[Long-Term Dependency]
  \label{def:long-term}
  A prediction task exhibits \emph{long-term dependency} when the evidence
  needed to make a correct prediction lies many time steps before the
  prediction point, so that the gap between relevant context and usage is
  large.
\end{definition}

The canonical classroom example from the slide deck (V01, 00:18:55) is
reproduced below.

\begin{example}[title=Karnataka/Delhi/Hindi Example (V01)]
  \textit{``I am from a South Indian family, specifically Karnataka, but
  settled in Delhi. I had spent most of my childhood and done schooling as
  well as college education in Delhi. Apart from Kannada, I can speak
  fluent \underline{\hspace{1.5em}} also.''}

  \medskip
  \noindent The missing word is ``Hindi''. The \emph{short-term} context
  (``Apart from Kannada, I can speak fluent'') tells us to expect a
  \emph{language name} but is ambiguous about \emph{which} language.
  To resolve the ambiguity we need the long-range context
  ``settled in Delhi'' --- many tokens earlier. A standard RNN with
  vanishing gradients cannot bridge this gap.
\end{example}

% =============================================================================
\section{Why Standard RNNs Fail on Long-Range Context}
\label{sec:rnn-failure}
% =============================================================================

\subsection{Backpropagation Through Time (BPTT)}
\label{subsec:bptt}

Training RNNs uses Backpropagation Through Time (BPTT), which unrolls the
network across all $T$ time steps and applies the chain rule. The gradient
of the loss $\mathcal{L}$ with respect to the hidden state at an early time
step $t$ passes through every intermediate step:

\begin{align}
  \label{eq:bptt-chain}
  \frac{\partial \mathcal{L}}{\partial h_t}
  &= \frac{\partial \mathcal{L}}{\partial h_T}
     \prod_{k=t+1}^{T}
     \frac{\partial h_k}{\partial h_{k-1}}.
\end{align}

Each factor in the product is

\begin{equation}
  \label{eq:jacobian-factor}
  \frac{\partial h_k}{\partial h_{k-1}}
  = W_{hh}^{\top} \operatorname{diag}\!\bigl(\tanh'(a_{k-1})\bigr),
\end{equation}

where $a_{k-1} = W_{hh}h_{k-2} + W_{xh}x_{k-1} + b_h$ is the
pre-activation. Since $|\tanh'(z)| \leq 1$ for all $z$, the spectral norm
of each factor is bounded by $\|W_{hh}\|_2$.

\subsection{The Vanishing Gradient Problem}
\label{subsec:vanishing}

\begin{theorem}[Gradient Magnitude Decay]
  \label{thm:vanishing}
  Let $\gamma = \|W_{hh}\|_2 \cdot \max_z|\tanh'(z)|$. If $\gamma < 1$,
  then
  \begin{equation}
    \label{eq:vanishing-bound}
    \left\|\frac{\partial \mathcal{L}}{\partial h_t}\right\|
    \;\leq\;
    \left\|\frac{\partial \mathcal{L}}{\partial h_T}\right\|
    \cdot \gamma^{T-t}.
  \end{equation}
  The gradient at step $t$ decays \emph{exponentially} with the distance
  $T - t$ from the output.
\end{theorem}

\begin{example}[title=Numerical Illustration (V04 scenario)]
  Consider the sentence ``I am from France \ldots\ I speak fluent
  \underline{\hspace{1em}}'' (slide 20/22 of \texttt{06-RNNTypesBackpropagation}).
  Suppose the sequence has $T = 6$ steps and $\gamma \approx 0.5$. The
  gradient arriving at the first token (``France'') is
  \[
    \|\nabla_{h_1}\| \;\leq\; \|\nabla_{h_6}\| \cdot (0.5)^{5}
    \;=\; \|\nabla_{h_6}\| \cdot 0.03125.
  \]
  That is, only $\approx 3\%$ of the gradient signal from the output
  reaches the earliest relevant token --- effectively zero for learning.
\end{example}

\cref{fig:rnn-fail} illustrates the gradient flow for this scenario.

\begin{figure}[ht]
  \centering
  % ── TikZ: RNN Long-Dependency Failure ─────────────────────────────────────
  \begin{tikzpicture}[
      node distance=1.4cm and 1.8cm,
      cell/.style={draw,rounded corners=3pt,minimum width=1.0cm,
                   minimum height=0.8cm,fill=orange!20,font=\small},
      word/.style={font=\small\itshape},
      arr/.style={-{Stealth[scale=0.8]},thick},
      garr/.style={-{Stealth[scale=0.8]},thick,red!70!black,dashed},
    ]
    \node[cell] (h1) {$h_1$};
    \node[cell,right=of h1] (h2) {$h_2$};
    \node[cell,right=of h2] (h3) {$h_3$};
    \node[cell,right=of h3] (h4) {$h_4$};
    \node[cell,right=of h4] (h5) {$h_5$};
    \node[cell,right=of h5] (h6) {$h_6$};

    % forward arrows
    \draw[arr] (h1)--(h2);
    \draw[arr] (h2)--(h3);
    \draw[arr] (h3)--(h4);
    \draw[arr] (h4)--(h5);
    \draw[arr] (h5)--(h6);

    % input words
    \node[word,below=0.5cm of h1] (w1) {France};
    \node[word,below=0.5cm of h2] (w2) {\ldots};
    \node[word,below=0.5cm of h3] (w3) {\ldots};
    \node[word,below=0.5cm of h4] (w4) {\ldots};
    \node[word,below=0.5cm of h5] (w5) {speak};
    \node[word,below=0.5cm of h6] (w6) {fluent};

    \draw[arr] (w1)--(h1);
    \draw[arr] (w2)--(h2);
    \draw[arr] (w3)--(h3);
    \draw[arr] (w4)--(h4);
    \draw[arr] (w5)--(h5);
    \draw[arr] (w6)--(h6);

    % output
    \node[above=0.7cm of h6,font=\small\bfseries] (yhat) {$\hat{y}$: ``\textit{French}''};
    \draw[arr] (h6)--(yhat);

    % backward gradient arrows (dashed red, fading)
    \draw[garr,line width=1.8pt] (h6) to[bend left=25]
          node[above,font=\tiny,red!80!black]{$\nabla$} (h5);
    \draw[garr,line width=1.3pt] (h5) to[bend left=25]
          node[above,font=\tiny,red!80!black]{$\approx 0.5\nabla$} (h4);
    \draw[garr,line width=0.9pt] (h4) to[bend left=25] (h3);
    \draw[garr,line width=0.5pt] (h3) to[bend left=25] (h2);
    \draw[garr,line width=0.2pt] (h2) to[bend left=25]
          node[above,font=\tiny,red!80!black]{$\approx 0$} (h1);

    % annotation
    \node[below=1.2cm of h1,text=red!70!black,font=\small\itshape,
          align=center] {Gradient nearly\\vanishes here};
  \end{tikzpicture}
  \caption{Gradient flow in a plain RNN processing a long-dependency
           sentence. The dashed red arrows show the gradient magnitude
           decaying from right (output) to left (early token ``France'').
           By the time the gradient reaches $h_1$, it is negligibly small,
           so the network cannot learn that ``France'' determines the
           answer.}
  \label{fig:rnn-fail}
\end{figure}

\subsection{Exploding Gradients}
\label{subsec:exploding}

The complementary failure mode occurs when $\gamma > 1$: the gradient grows
\emph{exponentially}, causing numerical overflow and training instability.
As \cref{eq:vanishing-bound} shows, the regime $\gamma < 1$ causes
vanishing and $\gamma > 1$ causes explosion --- both are pathological.

\begin{warningbox}[title=Vanishing vs Exploding Gradient]
  The plain RNN \cref{eq:rnn-update} is unstable in the BPTT sense:
  \begin{itemize}[nosep]
    \item If $\|W_{hh}\|_2 \cdot \max|\tanh'| < 1$: gradients
          \textbf{vanish} --- long-term dependencies cannot be learned.
    \item If $\|W_{hh}\|_2 \cdot \max|\tanh'| > 1$: gradients
          \textbf{explode} --- training diverges.
  \end{itemize}
  Gradient clipping (heuristic) can partially address explosion, but
  vanishing requires an architectural fix --- which is precisely what
  LSTM provides.
\end{warningbox}

\subsection{Short-Term vs Long-Term: A Practical Comparison}
\label{subsec:comparison}

\cref{tab:rnn-lstm-comparison} summarises the empirical performance
difference observed by the instructors (confirmed via a TA experiment on
sentiment classification):

\begin{table}[ht]
  \centering
  \caption{RNN vs LSTM accuracy (\%) on sentiment classification for
           varying sequence lengths. Longer sequences magnify the RNN
           performance gap (illustrative values from lecture).}
  \label{tab:rnn-lstm-comparison}
  \begin{tabular}{@{}lcc@{}}
    \toprule
    \textbf{Sequence Length} & \textbf{Plain RNN (\%)} & \textbf{LSTM (\%)} \\
    \midrule
    50  & 81.0 & 81.6 \\
    100 & 81.6 & 82.0 \\
    200 & $\downarrow$ (degraded) & 85.0 \\
    \midrule
    \multicolumn{3}{@{}l@{}}{\small As sequence length increases, RNN
    accuracy degrades gracefully then plateaus at a lower value.} \\
    \bottomrule
  \end{tabular}
\end{table}

\begin{keyinsight}[title=Task-Dependent Architecture Choice]
  It is \emph{not} always the case that LSTM outperforms RNN. For
  purely short-term tasks, LSTM adds computational overhead (roughly
  $4\times$ more parameters per cell) with negligible accuracy gain.
  LSTM is warranted when the task genuinely requires long-range context.
\end{keyinsight}

% =============================================================================
\section{The LSTM Architecture: Motivation and Design}
\label{sec:lstm-design}
% =============================================================================

\subsection{Design Rationale: Two Feedback Lines}
\label{subsec:two-lines}

The instructors motivated the LSTM design from first principles. The
insight is that a plain RNN has \emph{one} feedback line (the hidden state
$h_{t-1}$) which conflates short-term and long-term context:

\begin{quote}
``The first [idea] was we need two feedback parts. One, to keep track of the
short-term context information. Another to keep track of the long-term
context information.'' \hfill (Prof.\ Prasanna, 00:25)
\end{quote}

This motivates the dual-state design:

\begin{itemize}
  \item \textbf{Hidden state} $h_t$: the lower feedback line; carries
        short-term sequence knowledge (identical in spirit to a plain RNN
        hidden state).
  \item \textbf{Cell state} $C_t$: the upper feedback line; a
        \emph{memory highway} that preserves long-term context across
        potentially many time steps.
\end{itemize}

Between these two lines, \textbf{gate mechanisms} regulate the flow of
information: deciding what to write to $C_t$, what to erase from $C_t$,
and how much of $C_t$ to expose in $h_t$.

\subsection{The Cell State Highway}
\label{subsec:cell-state}

The cell state $C_t$ is the defining innovation. As Prof.\ Prasanna
explained:

\begin{quote}
``The cell state is the horizontal line running through the top of the
diagram. LSTM can remove or add information to the cell state, regulated by
structures called gates.''
\end{quote}

The key property is that updates to $C_t$ are \emph{additive}, not
multiplicative:

\begin{equation}
  \label{eq:cell-state-update}
  C_t = f_t \odot C_{t-1} + i_t \odot \tilde{C}_t,
\end{equation}

where $\odot$ denotes element-wise (Hadamard) multiplication, $f_t$ is
the forget-gate vector, $i_t$ is the input-gate vector, and $\tilde{C}_t$
is the candidate cell state. The gradient of $C_t$ with respect to
$C_{t-1}$ is simply $f_t$ --- a diagonal matrix with entries in $[0,1]$.
There is no repeated multiplication by $W_{hh}$ or $\tanh'$ along this
path, so gradients can flow unobstructed over long distances.

\begin{keyinsight}[title=Additive Updates Prevent Vanishing Gradients]
  In a plain RNN the gradient through time is a \emph{product} of Jacobians
  (\cref{eq:bptt-chain}), causing exponential decay. On the LSTM cell-state
  highway the gradient path reduces to a \emph{sum} modulated only by the
  forget gate. This is the fundamental mechanism by which LSTM overcomes the
  vanishing gradient problem.
\end{keyinsight}

% =============================================================================
\section{Gates in LSTM}
\label{sec:gates}
% =============================================================================

\subsection{The Gate Concept}
\label{subsec:gate-concept}

A \textbf{gate} in LSTM is a combination of a sigmoid neural network and a
pointwise multiplication (slide 37/63 of \texttt{05-DL-RNN}, V02 at
00:38:50):

\begin{equation}
  \label{eq:gate-generic}
  g_t = \sigma\!\bigl(W_g\,[h_{t-1},\,x_t] + b_g\bigr),
\end{equation}

where $\sigma(\cdot)$ is the sigmoid function, $[h_{t-1}, x_t]$ denotes
the concatenation of the previous hidden state and the current input, and
$W_g$, $b_g$ are learnable parameters specific to gate $g$. The output
$g_t \in [0,1]^n$ is then applied element-wise to the vector being gated:

\begin{equation}
  \label{eq:gate-apply}
  \text{gated output} = g_t \odot \text{input vector}.
\end{equation}

\begin{keyinsight}[title=Sigmoid as a Soft Gate]
  Unlike a digital logic gate (binary 0 or 1), an LSTM gate produces
  \emph{continuous} values in $[0,1]$ per dimension. A value of 0 means
  ``block this dimension completely''; a value of 1 means ``pass this
  dimension through unchanged''; intermediate values scale the
  information proportionally. This differentiability is what allows
  backpropagation to train the gate parameters.
\end{keyinsight}

\cref{fig:gate-diagram} shows the structure of a single LSTM gate.

\begin{figure}[ht]
  \centering
  \begin{tikzpicture}[
      node distance=1.2cm,
      sig/.style={draw,circle,minimum size=0.8cm,fill=yellow!30,font=\small},
      mul/.style={draw,circle,minimum size=0.8cm,fill=pink!50,font=\small},
      arr/.style={-{Stealth[scale=0.8]},thick},
      lbl/.style={font=\small},
    ]
    % input
    \node[lbl] (inp) {$[h_{t-1},\,x_t]$};
    % sigmoid block
    \node[sig,right=1.5cm of inp] (sig) {$\sigma$};
    % multiply
    \node[mul,right=2.0cm of sig] (mult) {$\times$};
    % info stream from above
    \node[lbl,above=1.2cm of mult] (stream) {information vector};
    % output
    \node[lbl,right=1.5cm of mult] (out) {gated output};

    \draw[arr] (inp)--(sig) node[midway,above,font=\tiny]{weight $W_g$};
    \draw[arr] (sig)--(mult) node[midway,above,font=\small]{$g_t \in [0,1]^n$};
    \draw[arr] (stream)--(mult);
    \draw[arr] (mult)--(out);

    % annotation boxes
    \node[draw,dashed,rounded corners,
          fit=(inp)(sig),inner sep=4pt,
          label=below:{\footnotesize sigmoid net}] {};
    \node[draw,dashed,rounded corners,
          fit=(mult),inner sep=6pt,
          label=below:{\footnotesize pointwise $\odot$}] {};
  \end{tikzpicture}
  \caption{Structure of a single LSTM gate (slide 37/63, V02). The sigmoid
           layer maps the concatenated input to a gating vector $g_t \in
           [0,1]^n$, which is then multiplied element-wise with the
           information vector to selectively allow or block information.}
  \label{fig:gate-diagram}
\end{figure}

\subsection{The Three Gates of LSTM}
\label{subsec:three-gates}

LSTM uses three gates, each with independent weight matrices, to regulate
the cell state $C_t$:

\subsubsection{Forget Gate $f_t$}
\label{subsubsec:forget}

The forget gate decides how much of the \emph{previous} cell state $C_{t-1}$
to retain:

\begin{equation}
  \label{eq:forget-gate}
  f_t = \sigma\!\bigl(W_f\,[h_{t-1},\,x_t] + b_f\bigr).
\end{equation}

When $f_t \approx \mathbf{0}$, the previous cell state is completely
erased (``forgotten''). When $f_t \approx \mathbf{1}$, it is fully
preserved. As the instructor stated:

\begin{quote}
``One means completely keep this dimension, zero means completely get rid of
this dimension. So that is the purpose of forgetting it.''
\hfill (Prof.\ Prasanna, 00:41)
\end{quote}

\subsubsection{Input Gate $i_t$ and Candidate Cell State $\tilde{C}_t$}
\label{subsubsec:input}

The input gate decides which components of the \emph{current} information
should be written to the cell state:

\begin{align}
  \label{eq:input-gate}
  i_t &= \sigma\!\bigl(W_i\,[h_{t-1},\,x_t] + b_i\bigr), \\
  \label{eq:candidate}
  \tilde{C}_t &= \tanh\!\bigl(W_C\,[h_{t-1},\,x_t] + b_C\bigr).
\end{align}

The candidate $\tilde{C}_t \in [-1,1]^n$ is the new information proposed
by the standard RNN operation ($\tanh$ of the weighted input). The gate
$i_t$ decides how much of this candidate actually enters the cell state.

\begin{remark}
  \label{rem:input-gate-naming}
  The instructors noted a common naming confusion: in Gated Recurrent
  Units (GRU) the analogous gate is called the \emph{update gate}, whereas
  in LSTM it is called the \emph{input gate}. Both serve the same
  conceptual purpose of writing new information into the memory.
\end{remark}

\subsubsection{Cell State Update}
\label{subsubsec:cell-update}

The cell state is updated by combining the forget and input gate outputs:

\begin{equation}
  \label{eq:cell-update}
  C_t = f_t \odot C_{t-1} + i_t \odot \tilde{C}_t.
\end{equation}

The additive structure of \cref{eq:cell-update} is what makes $C_t$ a
gradient highway: $\partial C_t / \partial C_{t-1} = f_t$, which is not
repeatedly multiplied across time steps.

\subsubsection{Output Gate $o_t$}
\label{subsubsec:output}

The output gate determines how much of the cell state to expose as the
hidden state $h_t$:

\begin{align}
  \label{eq:output-gate}
  o_t &= \sigma\!\bigl(W_o\,[h_{t-1},\,x_t] + b_o\bigr), \\
  \label{eq:hidden-update}
  h_t &= o_t \odot \tanh(C_t).
\end{align}

The $\tanh$ squashes $C_t$ into $[-1,1]$ before the gate is applied,
keeping the hidden state bounded.

\subsection{Summary of LSTM Equations}
\label{subsec:lstm-equations}

The complete set of LSTM equations for one time step is:

\begin{align}
  f_t          &= \sigma(W_f\,[h_{t-1},x_t] + b_f)
                 \label{eq:lstm-f} \\
  i_t          &= \sigma(W_i\,[h_{t-1},x_t] + b_i)
                 \label{eq:lstm-i} \\
  \tilde{C}_t  &= \tanh(W_C\,[h_{t-1},x_t] + b_C)
                 \label{eq:lstm-Ctil} \\
  C_t          &= f_t \odot C_{t-1} + i_t \odot \tilde{C}_t
                 \label{eq:lstm-C} \\
  o_t          &= \sigma(W_o\,[h_{t-1},x_t] + b_o)
                 \label{eq:lstm-o} \\
  h_t          &= o_t \odot \tanh(C_t)
                 \label{eq:lstm-h}
\end{align}

\cref{tab:lstm-params} catalogues the learnable parameter matrices.

\begin{table}[ht]
  \centering
  \caption{Learnable parameters in a single LSTM cell with hidden size $n$
           and input size $d$. Each gate has its own weight matrix and bias.}
  \label{tab:lstm-params}
  \begin{tabular}{@{}llll@{}}
    \toprule
    \textbf{Gate / State} & \textbf{Weight Matrix} & \textbf{Shape} & \textbf{Bias Shape} \\
    \midrule
    Forget gate $f_t$      & $W_f$ & $n \times (n+d)$ & $n$ \\
    Input gate $i_t$       & $W_i$ & $n \times (n+d)$ & $n$ \\
    Candidate $\tilde{C}_t$ & $W_C$ & $n \times (n+d)$ & $n$ \\
    Output gate $o_t$      & $W_o$ & $n \times (n+d)$ & $n$ \\
    \midrule
    \multicolumn{4}{@{}l@{}}{\small Total parameters: $4n(n+d) + 4n = 4n(n+d+1)$} \\
    \bottomrule
  \end{tabular}
\end{table}

% =============================================================================
\section{LSTM Cell Architecture: Diagram and Walk-Through}
\label{sec:lstm-cell}
% =============================================================================

\subsection{The Full LSTM Cell Diagram}
\label{subsec:lstm-cell-diagram}

\cref{fig:lstm-cell} shows the full LSTM recurrence cell in the style of
Colah's blog (referenced in slides 35/63 of \texttt{05-DL-RNN}, V03 and V05).

\begin{figure}[ht]
  \centering
  \begin{tikzpicture}[
      scale=0.95,
      node distance=0.9cm and 1.1cm,
      sig/.style={draw,rounded corners=2pt,minimum width=0.7cm,
                  minimum height=0.6cm,fill=yellow!40,font=\small},
      tnh/.style={draw,rounded corners=2pt,minimum width=0.7cm,
                  minimum height=0.6cm,fill=cyan!30,font=\small},
      op/.style={draw,circle,minimum size=0.55cm,fill=pink!50,font=\footnotesize},
      arr/.style={-{Stealth[scale=0.7]},thick},
      lbl/.style={font=\small},
    ]

    % ── Cell state (top horizontal line) ──────────────────────────────────
    \coordinate (CL) at (0,3.5);
    \coordinate (CR) at (10,3.5);

    % ── Forget gate ────────────────────────────────────────────────────────
    \node[sig]  (fg)  at (2.0,1.5) {$\sigma$};
    \node[op]   (fmul) at (2.0,3.5) {$\times$};

    % ── Input gate + candidate ─────────────────────────────────────────────
    \node[sig]  (ig)  at (4.5,1.5) {$\sigma$};
    \node[tnh]  (cand) at (6.0,1.5) {$\tanh$};
    \node[op]   (imul) at (5.2,3.5) {$\times$};
    \node[op]   (iadd) at (7.0,3.5) {$+$};

    % ── Output gate ────────────────────────────────────────────────────────
    \node[sig]  (og)  at (8.5,1.5) {$\sigma$};
    \node[tnh]  (otanh) at (8.5,3.5) {$\tanh$};
    \node[op]   (omul) at (9.3,2.5) {$\times$};

    % ── Cell state line ────────────────────────────────────────────────────
    \draw[arr,line width=1.2pt,blue!60!black]
          (CL) -- (fmul) -- (imul) -- (iadd) -- (otanh);

    % ── Forget gate connections ────────────────────────────────────────────
    \draw[arr] (fg) -- (fmul);

    % ── Input gate connections ─────────────────────────────────────────────
    \draw[arr] (ig)   -- (imul);
    \draw[arr] (cand) -- (imul);
    \draw[arr] (imul) -- (iadd);

    % ── Output gate connections ────────────────────────────────────────────
    \draw[arr] (iadd) -- (otanh);
    \draw[arr] (otanh) -- (omul);
    \draw[arr] (og)  -- (omul);

    % ── Input h_{t-1} and x_t (shared concatenated input) ─────────────────
    \coordinate (htminus1) at (0.5, 0.2);
    \node[lbl,left=0.1cm of htminus1] {$h_{t-1}$};
    \coordinate (xt)       at (0.5, -0.5);
    \node[lbl,left=0.1cm of xt] {$x_t$};

    % feed all gates
    \draw[arr,gray!70] (htminus1) -- (fg.south);
    \draw[arr,gray!70] (xt)       -- (fg.south);
    \draw[arr,gray!70] (htminus1) to[out=0,in=240] (ig.south);
    \draw[arr,gray!70] (xt)       to[out=0,in=270] (ig.south);
    \draw[arr,gray!70] (htminus1) to[out=0,in=240] (cand.south);
    \draw[arr,gray!70] (xt)       to[out=0,in=270] (cand.south);
    \draw[arr,gray!70] (htminus1) to[out=0,in=240] (og.south);
    \draw[arr,gray!70] (xt)       to[out=0,in=270] (og.south);

    % ── Outputs ────────────────────────────────────────────────────────────
    \node[lbl,right=0.3cm of omul] (ht) {$h_t$};
    \draw[arr,line width=1.2pt] (omul) -- (ht);

    \node[lbl,above=0.2cm of CR] {$C_t$};
    \node[lbl,above=0.2cm of CL] {$C_{t-1}$};

    % ── Labels ─────────────────────────────────────────────────────────────
    \node[below=0.1cm of fg,  font=\footnotesize,text=gray] {forget};
    \node[below=0.1cm of ig,  font=\footnotesize,text=gray] {input};
    \node[below=0.1cm of cand,font=\footnotesize,text=gray] {candidate};
    \node[below=0.1cm of og,  font=\footnotesize,text=gray] {output};

  \end{tikzpicture}
  \caption{LSTM recurrence cell (after Colah's blog, referenced in slides
           35/63 of \texttt{05-DL-RNN}, V03). The thick blue horizontal line
           is the \textbf{cell state} $C_t$ --- the long-term memory highway.
           Yellow $\sigma$ blocks are sigmoid gates; cyan $\tanh$ blocks
           produce bounded values. Pink circles denote element-wise
           operations ($\times$ = multiply, $+$ = add).}
  \label{fig:lstm-cell}
\end{figure}

\subsection{Step-by-Step Walk-Through}
\label{subsec:walkthrough}

The following algorithm formalises the forward pass through one LSTM cell,
as explained by the instructors.

\begin{algorithm}[ht]
  \DontPrintSemicolon
  \caption{LSTM Cell Forward Pass}
  \label{alg:lstm-forward}
  \KwIn{Current input $x_t \in \mathbb{R}^d$, previous hidden state
        $h_{t-1} \in \mathbb{R}^n$, previous cell state
        $C_{t-1} \in \mathbb{R}^n$}
  \KwOut{New hidden state $h_t \in \mathbb{R}^n$, new cell state
         $C_t \in \mathbb{R}^n$}
  \BlankLine
  \tcp{Step 1: Concatenate inputs}
  $z \leftarrow [h_{t-1};\, x_t] \in \mathbb{R}^{n+d}$\;
  \BlankLine
  \tcp{Step 2: Forget gate -- decide what to erase from $C_{t-1}$}
  $f_t \leftarrow \sigma(W_f z + b_f)$\;
  \BlankLine
  \tcp{Step 3: Input gate -- decide what new info to write}
  $i_t \leftarrow \sigma(W_i z + b_i)$\;
  \BlankLine
  \tcp{Step 4: Candidate cell state -- new information proposal}
  $\tilde{C}_t \leftarrow \tanh(W_C z + b_C)$\;
  \BlankLine
  \tcp{Step 5: Update cell state (additive -- gradient highway)}
  $C_t \leftarrow f_t \odot C_{t-1} + i_t \odot \tilde{C}_t$\;
  \BlankLine
  \tcp{Step 6: Output gate -- decide how much of $C_t$ to expose}
  $o_t \leftarrow \sigma(W_o z + b_o)$\;
  \BlankLine
  \tcp{Step 7: Compute new hidden state}
  $h_t \leftarrow o_t \odot \tanh(C_t)$\;
  \BlankLine
  \Return{$h_t,\; C_t$}
\end{algorithm}

\subsection{Intuition Behind Each Step}
\label{subsec:intuition}

\begin{keyinsight}[title=Why Tanh for the Candidate, Sigmoid for Gates?]
  The candidate $\tilde{C}_t$ uses $\tanh$ to produce values in $[-1,1]$,
  which can represent both \emph{increasing} and \emph{decreasing} a
  dimension of the cell state. The gates use $\sigma$ to produce values in
  $[0,1]$, acting as \emph{soft switches}. As the instructor noted, any
  function with output range $[0,1]$ could in principle serve as a gate;
  sigmoid is used because it is well-established and differentiable
  everywhere.
\end{keyinsight}

The key intuitions for each gate, as articulated in the lecture, are:

\begin{itemize}
  \item \textbf{Forget gate} $f_t$: ``I check whether this long-term
        information which I'm preserving from a long duration, whether it
        is relevant now or we can throw it out.''
  \item \textbf{Input gate} $i_t$ + candidate $\tilde{C}_t$: ``Whether
        the current state information should become part of the long-term
        context information --- that is decided by this input gate.''
  \item \textbf{Output gate} $o_t$: ``Whatever the output I have, I want
        to convert [it] into the output for the next sequence of
        information.''
\end{itemize}

% =============================================================================
\section{How LSTM Solves the Vanishing Gradient Problem}
\label{sec:solution}
% =============================================================================

\subsection{The Gradient Path Through the Cell State}
\label{subsec:grad-path}

Consider backpropagating through the cell state update
(\cref{eq:cell-update}). The partial derivative of $C_t$ with respect to
$C_{t-1}$ is:

\begin{equation}
  \label{eq:cell-grad}
  \frac{\partial C_t}{\partial C_{t-1}}
  = \operatorname{diag}(f_t).
\end{equation}

Crucially, this is a \emph{diagonal} matrix whose entries are the forget
gate values $f_{t,j} \in [0,1]$. When tracing gradients across $k$ steps:

\begin{equation}
  \label{eq:cell-grad-multi}
  \frac{\partial C_t}{\partial C_{t-k}}
  = \prod_{\tau=1}^{k} \operatorname{diag}(f_{t-\tau+1}).
\end{equation}

If the network learns to set $f_t \approx \mathbf{1}$ for the relevant
dimensions, then $\prod f_{t-\tau+1} \approx \mathbf{1}$ and the gradient
flows unattenuated. This is fundamentally different from the RNN case where
the product involves the full matrix $W_{hh}$ at every step.

\begin{lemma}[LSTM Gradient Preservation]
  \label{lem:lstm-gradient}
  If $f_{t,j} = 1$ for all $t$ and some dimension $j$, then dimension $j$
  of the gradient $\partial \mathcal{L}/\partial C_{t_0,j}$ equals
  dimension $j$ of $\partial \mathcal{L}/\partial C_{T,j}$ for all
  $t_0 \leq T$ --- the gradient is perfectly preserved over arbitrarily
  many steps in that dimension.
\end{lemma}

\begin{keyinsight}[title=Additive vs Multiplicative Updates]
  The essential contrast between RNN and LSTM gradient flow:
  \begin{itemize}[nosep]
    \item \textbf{RNN}: $\partial h_t / \partial h_{t-1}$
          involves $W_{hh}$ and $\tanh'$ --- multiplicative, prone to
          vanishing/exploding.
    \item \textbf{LSTM cell state}: $\partial C_t / \partial C_{t-1}$
          is just the forget gate $f_t$ --- the gate is \emph{learned}
          and can be kept near 1, allowing gradients to flow.
  \end{itemize}
  A student in the session insightfully observed: ``We are doing a
  summation operator here\ldots\ even if we have very small values, we are
  not directly multiplying, we are adding to it.''
\end{keyinsight}

\subsection{Artificial Memory Lifting}
\label{subsec:artificial-memory}

The instructors provided a powerful intuition for what LSTM does
architecturally:

\begin{quote}
``If I just forget [early information], it means I don't go by short-term
[RNN], because of the vanishing gradient problem --- it'll disappear. That's
why I'm \emph{artificially lifting up} and putting into the long-term context
information and taking it forward.'' \hfill (Prof.\ Prasanna, 01:08)
\end{quote}

That is, whenever a piece of information is deemed important
(input gate $i_t \approx 1$), it is \emph{explicitly copied} from the
short-term hidden state into the cell-state highway, where it can travel
any number of steps without decaying. This is the architectural realisation
of the two-feedback-line intuition.

\cref{tab:rnn-lstm-arch} compares the two architectures side-by-side.

\begin{table}[ht]
  \centering
  \caption{Architectural comparison between standard RNN and LSTM.}
  \label{tab:rnn-lstm-arch}
  \begin{tabularx}{\linewidth}{@{}lXX@{}}
    \toprule
    \textbf{Property}
      & \textbf{Standard RNN}
      & \textbf{LSTM} \\
    \midrule
    Feedback lines
      & 1 ($h_t$)
      & 2 ($h_t$ + $C_t$) \\
    Memory type
      & Short-term only
      & Short-term ($h_t$) and long-term ($C_t$) \\
    Gradient path
      & Multiplicative: $W_{hh}^k \cdot \tanh'^k$
      & Additive through $C_t$; multiplicative for $h_t$ \\
    Parameter count
      & $n(n+d)+n$
      & $4n(n+d+1)$ \\
    Handles long deps.\
      & No (gradient vanishes)
      & Yes (cell-state highway) \\
    Complexity
      & Low
      & High ($\approx 4\times$ RNN) \\
    Use when\ldots
      & Short-term task
      & Long-term task required \\
    \bottomrule
  \end{tabularx}
\end{table}

% =============================================================================
\section{Worked Example: Language Modelling with Long Context}
\label{sec:worked-example}
% =============================================================================

\subsection{Scenario 1: Short Dependency (RNN Succeeds)}
\label{subsec:scenario1}

\begin{example}[title=``Your mobile number is selected as \underline{\hspace{1.5em}}'']
  Context: ``your mobile number is selected as'' (6 tokens). \medskip

  \noindent\textbf{RNN processing:}
  \begin{align*}
    h_1 &= \tanh(W_{xh}\,\mathbf{e}_{\text{your}} + b_h), \\
    h_2 &= \tanh(W_{hh}h_1 + W_{xh}\,\mathbf{e}_{\text{mobile}} + b_h), \\
        &\;\vdots \\
    h_6 &= \tanh(W_{hh}h_5 + W_{xh}\,\mathbf{e}_{\text{as}} + b_h), \\
    \hat{y} &= \text{softmax}(W_{hy}h_6).
  \end{align*}
  The word ``selected as'' strongly predicts ``winner''. The RNN handles
  this well because all the relevant evidence is within 6 steps.
\end{example}

\subsection{Scenario 2: Long Dependency (RNN Fails, LSTM Succeeds)}
\label{subsec:scenario2}

\begin{example}[title=``I am from France \ldots\ I speak fluent \underline{\hspace{1.5em}}'']
  Context (paraphrased from slide 20/22): ``I am from France. Once I visited
  India\ldots\ I speak fluent \underline{\hspace{1em}}''. The evidence
  ``France'' is at $h_1$, but the prediction is at $h_6$. \medskip

  \noindent\textbf{RNN gradient reaching $h_1$:}
  \[
    \|\nabla_{h_1}\| \leq \|\nabla_{h_6}\| \cdot \gamma^5
    \approx 0 \quad \text{for } \gamma < 1.
  \]
  \textbf{LSTM cell state at $h_1$:} The LSTM input gate fires
  ($i_1 \approx 1$) for the ``France'' dimension, writing it into $C_1$.
  The forget gate maintains it ($f_t \approx 1$) across steps 2--5. At
  step 6, the hidden state $h_6$ reads from $C_6 \approx C_1$ via the
  output gate, correctly predicting ``French''.
\end{example}

% =============================================================================
\section{LSTM in Practice: Design Considerations}
\label{sec:practice}
% =============================================================================

\subsection{When to Use LSTM vs RNN}
\label{subsec:when-to-use}

A recurring theme in the lecture was that LSTM is not a universal
replacement for RNN. The choice depends on the task:

\begin{itemize}
  \item \textbf{Short-term only tasks} (e.g.\ next word prediction from a
        short phrase): Standard RNN is sufficient and computationally
        cheaper.
  \item \textbf{Long-term tasks} (e.g.\ paragraph-level language modelling,
        machine translation, speech recognition): LSTM is required.
  \item \textbf{Mixed or unknown}: Use LSTM as the safer default, accepting
        the $\approx 4\times$ computation overhead.
\end{itemize}

\begin{warningbox}[title=LSTM is Not Always Better]
  ``It is not a generic statement to say that LSTM always gives better
  performance than RNN. It depends on the task at your hand. If your task
  needs only modelling of the short-term context information, then you will
  not gain much by LSTM.'' \hfill (Prof.\ Prasanna, 01:03)
\end{warningbox}

\subsection{Multiple Stacked LSTM Layers}
\label{subsec:stacked}

In practice, LSTM cells are stacked: the output $h_t^{(\ell)}$ of layer
$\ell$ becomes the input $x_t^{(\ell+1)}$ of layer $\ell+1$:

\begin{equation}
  \label{eq:stacked}
  h_t^{(\ell)},\; C_t^{(\ell)}
  = \operatorname{LSTM}^{(\ell)}\!\bigl(h_t^{(\ell-1)},\,
    h_{t-1}^{(\ell)},\, C_{t-1}^{(\ell)}\bigr),
  \quad \ell = 1,\ldots,L.
\end{equation}

Lower layers tend to capture syntactic patterns (short-range), while
higher layers capture semantic patterns (long-range), naturally creating a
hierarchy of dependency lengths.

\subsection{Bidirectional LSTM}
\label{subsec:bilstm}

Many NLP tasks benefit from both left-to-right and right-to-left context.
A Bidirectional LSTM (BiLSTM) runs two LSTM cells in opposite directions
and concatenates their hidden states:

\begin{equation}
  \label{eq:bilstm}
  h_t^{\text{bi}} = \bigl[\overrightarrow{h}_t;\; \overleftarrow{h}_t\bigr],
\end{equation}

where $\overrightarrow{h}_t$ is the forward LSTM hidden state and
$\overleftarrow{h}_t$ is the backward LSTM hidden state. This is
a standard feature of modern sequence-labelling systems.

% =============================================================================
\section{Summary and Conclusion}
\label{sec:conclusion}
% =============================================================================

\subsection{Session Summary}
\label{subsec:session-summary}

Session~19 covered the motivation and introduction of the LSTM architecture
through the lens of the vanishing gradient problem in standard RNNs.
\cref{tab:session-summary} maps each key concept to the visual evidence
from the lecture slides.

\begin{table}[ht]
  \centering
  \caption{Session~19 concept map: key ideas, visual evidence, and
           mathematical expression.}
  \label{tab:session-summary}
  \begin{tabularx}{\linewidth}{@{}lXl@{}}
    \toprule
    \textbf{Concept}
      & \textbf{Key Insight}
      & \textbf{Slide / Equation} \\
    \midrule
    Short-term dependency
      & ``Clouds are in the sky'' -- local context suffices
      & V01, \cref{def:short-term} \\
    Long-term dependency
      & Karnataka/Delhi/Hindi -- context spans paragraphs
      & V01, \cref{def:long-term} \\
    Vanishing gradient
      & $\gamma^{T-t} \to 0$ exponentially for $\gamma < 1$
      & V04, \cref{thm:vanishing} \\
    Gate concept
      & $\sigma \in [0,1]$ times information vector
      & V02, \cref{eq:gate-generic} \\
    Cell-state highway
      & Additive update; gradient $= f_t$
      & V03/V05, \cref{eq:cell-update} \\
    Forget gate
      & Erases stale long-term context
      & \cref{eq:forget-gate} \\
    Input gate
      & Writes new info to cell state
      & \cref{eq:input-gate,eq:candidate} \\
    Output gate
      & Exposes relevant cell state as $h_t$
      & \cref{eq:output-gate,eq:hidden-update} \\
    \bottomrule
  \end{tabularx}
\end{table}

\subsection{Key Takeaways}
\label{subsec:takeaways}

\begin{keyinsight}[title=The Three Core Ideas of Session 19]
  \begin{enumerate}[nosep,label=\textbf{\arabic*.}]
    \item \textbf{Why RNNs fail on long sequences.} Backpropagation
          through time multiplies Jacobians at every step; when the
          spectral radius is below 1, gradients vanish exponentially,
          erasing early context.
    \item \textbf{The LSTM fix: cell-state highway.} A second feedback
          line $C_t$ with additive (not multiplicative) updates provides a
          gradient-preserving memory highway across unlimited time steps.
    \item \textbf{Gates as soft switches.} Three sigmoid networks
          (forget, input, output) learn, from data, when to erase, write,
          and read the cell state --- enabling the network to maintain or
          discard long-term context as required by the task.
  \end{enumerate}
\end{keyinsight}

\subsection{Looking Ahead}
\label{subsec:looking-ahead}

Session~20 will perform a full numerical walkthrough of all three LSTM
gates for a concrete example, making the abstract equations in
\cref{eq:lstm-f,eq:lstm-i,eq:lstm-Ctil,eq:lstm-C,eq:lstm-o,eq:lstm-h}
tangible. The session will also introduce the Gated Recurrent Unit (GRU),
which simplifies the LSTM design to two gates while retaining most of the
long-term memory capability.

\medskip
\noindent\textbf{Acknowledgements.} The LSTM cell diagram in this document
is adapted from C.\ Olah, ``Understanding LSTM Networks,''
\textit{colah.github.io}, 2015. The slide content is from
\texttt{05-DL-RNN} and \texttt{06-RNNTypesBackpropagation},
IIITDwd, January 2026.

\clearpage

%% ============================================================
%% Session 20: LSTM Numerical Walkthrough and Next-Word Prediction Hands-On
%% ============================================================
\part{Session 20: LSTM Numerical Walkthrough and Next-Word Prediction Hands-On}

%% ============================================================
%%  TITLE PAGE
%% ============================================================
\begin{titlepage}
  \centering
  \vspace*{2cm}
  {\LARGE\bfseries Deep Learning: Course Notes\par}
  \vspace{0.8cm}
  \rule{\linewidth}{0.5pt}\\[0.4cm]
  {\huge\bfseries Session 20\\[0.3cm]
   LSTM Numerical Walkthrough\\[0.2cm]
   \& Next-Word Prediction Hands-On\par}
  \rule{\linewidth}{0.5pt}\\[1.5cm]

  \begin{tabular}{rl}
    \textbf{Session Number:} & 20 \\[4pt]
    \textbf{Date:}           & January 31, 2026 \\[4pt]
    \textbf{Instructors:}    & Prof.\ S.\ R.\ M.\ Prasanna \\
                             & Dr.\ Sunil Saumya \\[4pt]
    \textbf{Institute:}      & IIIT Dharwad \\[4pt]
    \textbf{Slide Deck:}     & \texttt{05-DL-RNN} (slides 35, 39, 42 of 75) \\[4pt]
    \textbf{Video:}          & \url{https://vimeo.com/1160427634} \\[4pt]
    \textbf{Duration:}       & $\approx 1$ hr 41 min \\
  \end{tabular}

  \vfill
  {\large\itshape Topics Covered}\\[0.4cm]
  \begin{itemize}[leftmargin=4cm,rightmargin=4cm]
    \item LSTM recurrence cell recap (Colah's blog style diagram)
    \item Forget Gate -- detailed numerical walkthrough
    \item Input Gate -- sigmoid and tanh sub-operations
    \item Output Gate -- filtered memory exposure
    \item Cell-state update equation
    \item GRU comparison (brief)
    \item Hands-on: next-word prediction with PyTorch LSTM
    \item Pre-processing pipeline: tokenisation, vocabulary, sliding-window sequences, padding
    \item \texttt{LSTMModel} class implementation
  \end{itemize}
  \vfill
\end{titlepage}

\newpage

%% ============================================================
\section{Introduction}
\label{sec:intro}
%% ============================================================

Session 20 of the deep learning course is devoted to two complementary goals.
The first half provides a careful, \emph{numerical} walkthrough of how the three
gates of a Long Short-Term Memory (LSTM) cell operate, working with explicit
4-dimensional vectors so that every dot-product, sigmoid evaluation, and
element-wise multiplication can be traced by hand.
The second half transitions to a fully working hands-on session in which an
LSTM is trained to perform \emph{next-word prediction} on a short text corpus,
walking through every stage of the PyTorch pipeline: tokenisation, vocabulary
construction, self-supervised sequence creation, padding, custom dataset and
data loader, model definition, and training.

\begin{keyinsight}
  The central motivation for LSTM is the \emph{vanishing gradient} problem of
  vanilla RNNs.  A plain RNN uses a \emph{single} hidden-state vector that must
  simultaneously encode short-term context (the current word) and long-term
  context (words seen many steps ago).  As back-propagation is applied through
  many time steps, gradients shrink exponentially, making it practically
  impossible to learn long-range dependencies.
  LSTM introduces a dedicated \textbf{cell state} $C_t$ that acts as a
  long-term memory highway, updated only through controlled additive operations
  that let gradients flow almost unchanged across many time steps.
\end{keyinsight}

The session is taught by Dr.\ Sunil Saumya (theoretical portion) and references
Prof.\ S.\ R.\ M.\ Prasanna's slides from the \texttt{05-DL-RNN} deck.
The visual material is drawn from Christopher Olah's widely cited blog post
``Understanding LSTM Networks'' (2015), which provides the iconic cell-diagram
used throughout the lecture.

%% ============================================================
\section{Background and Prerequisites}
\label{sec:background}
%% ============================================================

\subsection{Vanilla RNN and Its Limitations}
\label{subsec:rnn-limits}

A standard (Elman) RNN maintains a single hidden state $h_t$ updated as

\begin{equation}
  h_t = \tanh\!\bigl(W_h h_{t-1} + W_x x_t + b\bigr),
  \label{eq:rnn-update}
\end{equation}

where $x_t \in \R^d$ is the input at time step $t$, $W_h$ and $W_x$ are
learned weight matrices, and $b$ is a bias vector.
The same vector $h_t$ serves as both:
\begin{itemize}
  \item the \textbf{short-term} representation passed forward to the next cell, and
  \item the \textbf{long-term} carrier of all contextual information seen so far.
\end{itemize}
Because $h_t$ is overwritten at every step, information from distant time steps
(e.g.\ $x_{100}$) is diluted by the time the network processes $x_1$ in
reverse during BPTT, resulting in near-zero gradients --- the vanishing
gradient problem.

\subsection{Word Representations}
\label{subsec:word-rep}

Before feeding text into any recurrent network, words must be converted to
vectors.  Two representations are used in this session:

\begin{description}[style=nextline]
  \item[One-hot vectors] A vocabulary of size $V$ maps each word to a binary
    vector $v \in \{0,1\}^V$ with exactly one non-zero entry.  This is
    sparse, high-dimensional, and carries no semantic information.
  \item[Dense embeddings] An \texttt{nn.Embedding} layer learns a weight matrix
    $E \in \R^{V \times d_e}$; look-up of word index $i$ returns the $i$-th
    row, producing a dense $d_e$-dimensional vector.  This dramatically reduces
    dimensionality and allows the network to learn semantic similarity.
\end{description}

\begin{remark}
  In the running example used during the lecture, the vocabulary contains
  $V = 4$ words (\emph{I}, \emph{love}, \emph{deep}, \emph{learning}) and the
  embedding dimension is set to $d_e = 4$, giving equal embedding and hidden
  dimensions for simplicity.  In the later hands-on, $V = 289$, $d_e = 100$.
\end{remark}

\subsection{Tensor Dimensions for LSTM Input}
\label{subsec:tensor-dims}

With \texttt{batch\_first=True} (PyTorch convention used in this session),
the input tensor has shape

\[
  \text{input} \;\in\; \R^{B \times T \times d_e},
\]

where $B$ is the batch size, $T$ is the sequence length (number of time steps),
and $d_e$ is the embedding dimension.  For the minimal toy example
($B=1$, $T=4$, $d_e=4$) the shape is $1 \times 4 \times 4$.

%% ============================================================
\section{LSTM Cell Architecture}
\label{sec:lstm-arch}
%% ============================================================

\subsection{The Two-Stream Design}
\label{subsec:two-stream}

The defining architectural choice of an LSTM is the separation of internal
memory into two distinct vectors:

\begin{itemize}
  \item \textbf{Cell state} $C_t \in \R^H$ --- the long-term memory highway.
    It flows from left to right across the unrolled cell with only two
    point-wise operations (one multiplication and one addition), so gradients
    can propagate almost unchanged over many steps.
  \item \textbf{Hidden state} $h_t \in \R^H$ --- the short-term, context-filtered
    output.  It is derived from $C_t$ at each step and is what the cell exposes
    to the next time step and to downstream layers.
\end{itemize}

\begin{keyinsight}[title=Memory Storage vs.\ Memory Exposure]
  The instructor uses an evocative analogy: think of $C_t$ as everything stored
  in your brain --- childhood memories, professional knowledge, personal plans
  --- all of which is available but not all simultaneously relevant.  When
  someone asks you about LSTM, you draw on only the relevant subset and
  express it as your answer; that filtered response is $h_t$.
  LSTM explicitly \emph{separates memory storage} ($C_t$) from
  \emph{memory exposure} ($h_t$).
\end{keyinsight}

\subsection{TikZ Diagram: LSTM Recurrence Cell}
\label{subsec:lstm-diagram}

\cref{fig:lstm-cell} recreates the Colah's-blog style cell diagram shown on
slide 35/75 of the lecture deck.

\begin{figure}[ht]
\centering
\begin{tikzpicture}[
    >=Stealth,
    node distance=1.3cm and 1.6cm,
    gate/.style={draw, rounded corners=4pt, fill=yellow!30,
                 minimum width=1.1cm, minimum height=0.75cm, font=\small},
    op/.style={draw, circle, fill=pink!50, minimum size=0.55cm, font=\small\bfseries},
    vec/.style={->, thick},
    label/.style={font=\small\itshape}
  ]

  %% Cell state line (top horizontal highway)
  \node (cin)  [left=2.5cm of {(0,1.8)}] {};
  \node (cout) [right=5.8cm of cin]       {};
  \draw[vec, line width=1.5pt, blue!60] (cin) -- node[above, label]{$C_{t-1}$} ++(2.0,0)
       node (cmul) [op] {$\times$};
  \draw[vec, line width=1.5pt, blue!60] (cmul) -- ++(1.5,0)
       node (cadd) [op] {$+$};
  \draw[vec, line width=1.5pt, blue!60] (cadd) -- node[above, label]{$C_t$} ++(2.2,0);

  %% Forget gate
  \node (fg) [gate, below=1.4cm of cmul] {$\sigma_f$};
  \draw[vec] (fg) -- node[right, label, xshift=2pt]{$f_t$} (cmul);

  %% Input gate (sigmoid part)
  \node (ig) [gate, below=1.4cm of cadd, xshift=-0.7cm] {$\sigma_i$};
  %% Candidate (tanh part)
  \node (tg) [gate, right=0.5cm of ig] {$\tanh$};
  %% Multiply i_t and ~C_t
  \node (imul) [op, above=0.7cm of ig, xshift=0.5cm] {$\times$};
  \draw[vec] (ig)  -- (imul);
  \draw[vec] (tg)  -- (imul);
  \draw[vec] (imul) -- (cadd);

  %% Output gate
  \node (tanh2) [gate, right=1.2cm of cadd, yshift=-1.0cm] {$\tanh$};
  \node (og)    [gate, below=0.6cm of tanh2, xshift=-1.0cm] {$\sigma_o$};
  \node (omul)  [op,   right=0.5cm of og, yshift=0.5cm] {$\times$};
  \draw[vec] (cadd)  |- (tanh2);
  \draw[vec] (tanh2) -- (omul);
  \draw[vec] (og)    -- (omul);
  \draw[vec] (omul)  -- node[right, label]{$h_t$} ++(0,1.2);

  %% Input lines for x_t and h_{t-1}
  \node (xt)   [below=2.5cm of fg, xshift=-0.5cm]  {$x_t$};
  \node (htm1) [left=0.3cm of fg, yshift=-1.2cm]   {$h_{t-1}$};
  \draw[vec, dashed] (xt)   -| (fg);
  \draw[vec, dashed] (htm1) |- (fg);
  \draw[vec, dashed] (xt)   -| (ig);
  \draw[vec, dashed] (xt)   -| (tg);
  \draw[vec, dashed] (xt)   -| (og);

\end{tikzpicture}
\caption{LSTM recurrence cell (Colah's blog style, slide 35/75).
  Blue: cell-state highway $C_t$.
  Yellow rectangles: learned single-layer sub-networks (gates).
  Pink circles: point-wise operations ($\times$ = element-wise multiply,
  $+$ = element-wise add).
  The four internal networks are: forget gate ($\sigma_f$), input gate
  ($\sigma_i$), candidate generator ($\tanh$), and output gate ($\sigma_o$).}
\label{fig:lstm-cell}
\end{figure}

%% ============================================================
\section{LSTM Gate Mathematics and Numerical Walkthrough}
\label{sec:gates}
%% ============================================================

This section presents the complete forward-pass equations for one LSTM time
step, followed by the 4-dimensional numerical example worked through by the
instructor.

\subsection{Notation and Setup}
\label{subsec:notation}

Throughout, the LSTM has $H$ hidden units (neurons per gate).  At time step
$t$ the following vectors are available:

\begin{center}
\begin{tabular}{lll}
\toprule
\textbf{Symbol} & \textbf{Shape} & \textbf{Description} \\
\midrule
$x_t$     & $\R^d$  & Input embedding at time $t$ \\
$h_{t-1}$ & $\R^H$  & Hidden state from previous step \\
$C_{t-1}$ & $\R^H$  & Cell state from previous step \\
$f_t$     & $\R^H$  & Forget gate output \\
$i_t$     & $\R^H$  & Input gate output \\
$\tilde{C}_t$ & $\R^H$ & Candidate cell values \\
$o_t$     & $\R^H$  & Output gate \\
$C_t$     & $\R^H$  & Updated cell state \\
$h_t$     & $\R^H$  & Updated hidden state \\
\bottomrule
\end{tabular}
\end{center}

In the lecture example, $d = H = 4$ (4-dimensional vectors throughout).

\subsection{Forget Gate}
\label{subsec:forget-gate}

\begin{definition}[Forget Gate]
  The forget gate $f_t$ is a vector in $[0,1]^H$ computed by a single-layer
  neural network with sigmoid activation.  It decides, for each dimension of
  the cell state, what fraction of $C_{t-1}$ to \emph{retain}.
\end{definition}

\begin{align}
  z_f &= W_{xf}\, x_t + W_{hf}\, h_{t-1} + b_f, \label{eq:zf} \\
  f_t &= \sigma(z_f), \label{eq:ft}
\end{align}

where $W_{xf} \in \R^{H \times d}$, $W_{hf} \in \R^{H \times H}$,
$b_f \in \R^H$.

\textbf{Interpretation.}
Each element $f_t[k] \in [0,1]$:
\begin{itemize}
  \item $f_t[k] \approx 1$: keep almost all of $C_{t-1}[k]$.
  \item $f_t[k] \approx 0$: erase almost all of $C_{t-1}[k]$.
  \item $f_t[k] \in (0,1)$: partial retention (soft gating).
\end{itemize}

The forget-gated cell state is then

\begin{equation}
  C_{t-1}^{\text{forget}} = f_t \odot C_{t-1}.
  \label{eq:forget-applied}
\end{equation}

\subsubsection{Numerical Example (Forget Gate)}
\label{subsubsec:forget-numerical}

\begin{example}[title=Forget Gate -- 4-D Numerical Example (slide 42/75)]
Given:
\[
  x_t = \begin{bmatrix} 0.2 \\ 0.6 \\ 0.1 \\ 0.9 \end{bmatrix},\quad
  h_{t-1} = \begin{bmatrix} 0.5 \\ 0.1 \\ 0.3 \\ 0.7 \end{bmatrix},\quad
  C_{t-1} = \begin{bmatrix} 0.9 \\ 0.2 \\ 0.8 \\ 0.4 \end{bmatrix}.
\]

After applying \cref{eq:zf} with some fixed $W_{xf},W_{hf},b_f$:
\[
  z_f = \begin{bmatrix} 1.2 \\ -0.5 \\ 0.3 \\ 2.0 \end{bmatrix}.
\]

Applying sigmoid element-wise:
\[
  f_t = \sigma(z_f) = \begin{bmatrix}
    \sigma(1.2) \\ \sigma(-0.5) \\ \sigma(0.3) \\ \sigma(2.0)
  \end{bmatrix}
  = \begin{bmatrix} 0.77 \\ 0.38 \\ 0.57 \\ 0.88 \end{bmatrix}.
\]

Element-wise multiplication with $C_{t-1}$ (\cref{eq:forget-applied}):
\[
  f_t \odot C_{t-1}
  = \begin{bmatrix} 0.77 \\ 0.38 \\ 0.57 \\ 0.88 \end{bmatrix}
  \odot
  \begin{bmatrix} 0.9 \\ 0.2 \\ 0.8 \\ 0.4 \end{bmatrix}
  = \begin{bmatrix} 0.69 \\ 0.08 \\ 0.46 \\ 0.35 \end{bmatrix}.
\]

\textbf{Reading off the result:}
\begin{itemize}
  \item Dimension 0: 77\% of 0.9 retained $\to$ 0.69 (moderate forgetting).
  \item Dimension 1: only 38\% of 0.2 retained $\to$ 0.08 (mostly forgotten).
  \item Dimension 2: 57\% of 0.8 retained $\to$ 0.46 (partial forgetting).
  \item Dimension 3: 88\% of 0.4 retained $\to$ 0.35 (mostly kept).
\end{itemize}
\end{example}

\subsection{Input Gate and Candidate Cell State}
\label{subsec:input-gate}

The input gate has \emph{two} sub-components operating in parallel:

\begin{align}
  i_t         &= \sigma\!\bigl(W_{xi}\,x_t + W_{hi}\,h_{t-1} + b_i\bigr),
                \label{eq:it} \\
  \tilde{C}_t &= \tanh\!\bigl(W_{xC}\,x_t + W_{hC}\,h_{t-1} + b_C\bigr).
                \label{eq:Ct-cand}
\end{align}

\begin{description}
  \item[$\tilde{C}_t$ (candidate cell).] The $\tanh$ layer generates new
    \emph{candidate values} in $[-1,1]^H$ --- what new information \emph{could}
    be added to memory based on the current context.
  \item[$i_t$ (input gate).] The sigmoid layer decides, for each dimension,
    \emph{how much} of the candidate to actually write into memory.
\end{description}

\begin{keyinsight}[title=Analogy for Input Gate]
  Imagine learning that ``deep learning is a subfield of machine learning.''
  The candidate $\tilde{C}_t$ represents \emph{all} the new information
  (everything about deep learning).  The input gate $i_t$ decides that you
  should add only \emph{part} of it --- perhaps already knowing machine
  learning means the overlap can be discarded.
  So the actual addition to memory is $i_t \odot \tilde{C}_t$: only
  the fraction $i_t[k]$ of each candidate dimension is written.
\end{keyinsight}

\subsection{Cell State Update}
\label{subsec:cell-update}

The new cell state combines the forget-gated old memory with the input-gated
new information:

\begin{equation}
  \boxed{C_t = f_t \odot C_{t-1} \;+\; i_t \odot \tilde{C}_t.}
  \label{eq:cell-update}
\end{equation}

This is the core LSTM update equation.  Crucially, it is purely
\emph{additive}: information is subtracted (via multiplication by $f_t < 1$)
or added (via $i_t \odot \tilde{C}_t$), but there is no tanh squashing of the
highway itself, which is what allows gradients to flow far back in time.

\subsection{Output Gate and Hidden State}
\label{subsec:output-gate}

Finally, the output gate decides what part of the (now-updated) cell state
to expose as the hidden state:

\begin{align}
  o_t &= \sigma\!\bigl(W_{xo}\,x_t + W_{ho}\,h_{t-1} + b_o\bigr),
         \label{eq:ot} \\
  h_t &= o_t \odot \tanh(C_t).
         \label{eq:ht}
\end{align}

The $\tanh(C_t)$ squashes cell-state values into $[-1,1]$, and $o_t$ acts as
a soft mask that selects which dimensions are relevant for the current output.

\begin{remark}
  A student question during the lecture highlighted that all three sigmoid
  outputs ($f_t$, $i_t$, $o_t$) share the same functional form
  (linear combination followed by sigmoid) and even the same inputs
  ($x_t$ and $h_{t-1}$).  What distinguishes them are:
  (1) separate weight matrices ($W_{*f}$, $W_{*i}$, $W_{*o}$) that are
  independently learned; and
  (2) the object each gate is \emph{applied to}: $f_t$ acts on $C_{t-1}$,
  $i_t$ acts on the candidate $\tilde{C}_t$, and $o_t$ acts on $\tanh(C_t)$.
\end{remark}

\subsection{Complete LSTM Forward-Pass Summary}
\label{subsec:full-pass}

Collecting \cref{eq:ft,eq:it,eq:Ct-cand,eq:cell-update,eq:ot,eq:ht}:

\begin{align}
  f_t         &= \sigma(W_{xf}x_t + W_{hf}h_{t-1} + b_f) \label{eq:full-ft} \\
  i_t         &= \sigma(W_{xi}x_t + W_{hi}h_{t-1} + b_i) \label{eq:full-it} \\
  \tilde{C}_t &= \tanh(W_{xC}x_t + W_{hC}h_{t-1} + b_C) \label{eq:full-Ct} \\
  C_t         &= f_t \odot C_{t-1} + i_t \odot \tilde{C}_t \label{eq:full-cell} \\
  o_t         &= \sigma(W_{xo}x_t + W_{ho}h_{t-1} + b_o) \label{eq:full-ot} \\
  h_t         &= o_t \odot \tanh(C_t) \label{eq:full-ht}
\end{align}

\subsection{Dimensional Consistency Constraint}
\label{subsec:dim-constraint}

A key point stressed by the instructor: all gate outputs, the cell state,
and the hidden state must share the same dimension $H$.  This is enforced by
the number of neurons in each gate sub-network.

\begin{warningbox}[title=Dimension Mismatch Will Break the Forward Pass]
  If you set the number of neurons to $H' \neq H$, the gate outputs will be
  $\R^{H'}$-vectors while $C_{t-1} \in \R^H$.  The element-wise operations
  in \cref{eq:forget-applied,eq:cell-update} require identical shapes, so
  PyTorch will raise a dimension mismatch error.  PyTorch's \texttt{nn.LSTM}
  enforces this automatically via its \texttt{hidden\_size} parameter.
\end{warningbox}

\subsection{Gate Parameter Count}
\label{subsec:param-count}

For a single LSTM layer with input dimension $d$ and hidden size $H$,
the total number of trainable parameters is:

\begin{equation}
  \text{Params}_{\text{LSTM}} = 4 \times \bigl(H \cdot d + H \cdot H + H\bigr)
  = 4H(d + H + 1).
  \label{eq:param-count}
\end{equation}

The factor of 4 comes from the four sub-networks: forget gate, input gate,
candidate generator, and output gate.

\begin{example}[title=Parameter Count for the Hands-On Model]
  With $d = 100$ (embedding dimension) and $H = 150$ (hidden size):
  \[
    \text{Params}_{\text{LSTM}} = 4 \times 150 \times (100 + 150 + 1)
    = 600 \times 251 = 150\,600.
  \]
\end{example}

\newpage
\subsection{LSTM Algorithm}
\label{subsec:algorithm}

\begin{algorithm}[H]
\SetAlgoLined
\linesnumbered
\DontPrintSemicolon
\caption{LSTM Single Time-Step Forward Pass}
\label{alg:lstm-forward}
\KwIn{$x_t \in \R^d$, $h_{t-1} \in \R^H$, $C_{t-1} \in \R^H$,
      weight matrices $W_{xf}, W_{hf}, W_{xi}, W_{hi},
      W_{xC}, W_{hC}, W_{xo}, W_{ho}$,
      biases $b_f, b_i, b_C, b_o$}
\KwOut{$h_t \in \R^H$, $C_t \in \R^H$}

\BlankLine
\tcp{Step 1 -- Forget Gate}
$z_f \leftarrow W_{xf}\,x_t + W_{hf}\,h_{t-1} + b_f$\;
$f_t \leftarrow \sigma(z_f)$\;

\BlankLine
\tcp{Step 2 -- Input Gate (two sub-operations)}
$z_i \leftarrow W_{xi}\,x_t + W_{hi}\,h_{t-1} + b_i$\;
$i_t \leftarrow \sigma(z_i)$\;
$z_C \leftarrow W_{xC}\,x_t + W_{hC}\,h_{t-1} + b_C$\;
$\tilde{C}_t \leftarrow \tanh(z_C)$\;

\BlankLine
\tcp{Step 3 -- Cell State Update}
$C_t \leftarrow f_t \odot C_{t-1} + i_t \odot \tilde{C}_t$\;

\BlankLine
\tcp{Step 4 -- Output Gate}
$z_o \leftarrow W_{xo}\,x_t + W_{ho}\,h_{t-1} + b_o$\;
$o_t \leftarrow \sigma(z_o)$\;
$h_t \leftarrow o_t \odot \tanh(C_t)$\;

\BlankLine
\Return{$h_t,\; C_t$}\;
\end{algorithm}

%% ============================================================
\section{LSTM vs.\ GRU: A Brief Comparison}
\label{sec:gru}
%% ============================================================

The instructor briefly previewed the Gated Recurrent Unit (GRU), which will
be covered in detail in Session 21.  \cref{tab:lstm-gru} summarises the main
architectural differences.

\begin{table}[ht]
\centering
\caption{Comparison of LSTM and GRU architectures.}
\label{tab:lstm-gru}
\begin{tabularx}{\textwidth}{lXX}
\toprule
\textbf{Property}           & \textbf{LSTM}                              & \textbf{GRU} \\
\midrule
Memory vectors              & Two: $C_t$ (long-term) and $h_t$ (short-term) & One: $h_t$ only \\
Gates                       & Three gates + candidate: forget $f_t$, input $i_t$, output $o_t$, candidate $\tilde{C}_t$ & Two gates: reset $r_t$, update $z_t$ \\
Parameter count (input $d$, hidden $H$) & $4H(d+H+1)$ & $3H(d+H+1)$ \\
Cell state highway          & Explicit $C_t$ line                        & Merged into $h_t$ \\
Complexity                  & Higher                                     & Lower \\
Typical performance         & Strong on long sequences                   & Comparable; sometimes faster to train \\
Recommendation              & When sequence length is very long          & When computational efficiency matters \\
\bottomrule
\end{tabularx}
\end{table}

\begin{keyinsight}[title=When to Choose LSTM vs.\ GRU]
  The instructor's advice: ``If you care about efficiency and complexity,
  prefer GRU.  If you care about performance, you can choose either -- most of
  the time the difference is marginal.''  GRU achieves a similar effect to
  LSTM's forget/input split through its update gate, but uses a single hidden
  vector instead of the separate $C_t$ and $h_t$, reducing parameter count by
  approximately 25\%.
\end{keyinsight}

%% ============================================================
\section{Hands-On: Next-Word Prediction with PyTorch LSTM}
\label{sec:handson}
%% ============================================================

The second half of the session builds a complete next-word prediction
(language model) pipeline.  The corpus is a single paragraph of text
(approx.\ 78 sentences, 289 unique tokens) about a data science mentorship
programme, kept deliberately short to fit within a single session.

\subsection{Pipeline Overview}
\label{subsec:pipeline}

\cref{fig:pipeline} shows the end-to-end flow of data through the system.

\begin{figure}[ht]
\centering
\begin{tikzpicture}[
    >=Stealth,
    node distance=0.5cm and 0.1cm,
    box/.style={draw, rounded corners=4pt, fill=gray!10,
                minimum width=2.8cm, minimum height=0.9cm,
                text centered, font=\small},
    arr/.style={->, thick}
  ]
  \node (raw)   [box]                          {Raw Text\\Corpus};
  \node (lower) [box, right=1.1cm of raw]      {Lowercase +\\Tokenise (NLTK)};
  \node (vocab) [box, right=1.1cm of lower]    {Build\\Vocabulary};
  \node (idx)   [box, below=0.9cm of vocab]    {Text $\to$\\Index Sequences};
  \node (slide) [box, left=1.1cm of idx]       {Sliding-Window\\Self-Supervised};
  \node (pad)   [box, left=1.1cm of slide]     {Pad to\\$\max\_len=62$};
  \node (split) [box, below=0.9cm of pad]      {Split $X$, $y$\\(input, label)};
  \node (data)  [box, right=1.1cm of split]    {Custom Dataset\\+ DataLoader};
  \node (model) [box, right=1.1cm of data]     {LSTMModel\\(Emb+LSTM+FC)};
  \node (train) [box, right=1.1cm of model]    {Train\\(CE Loss + Adam)};

  \draw[arr] (raw)   -- (lower);
  \draw[arr] (lower) -- (vocab);
  \draw[arr] (vocab) -- (idx);
  \draw[arr] (idx)   -- (slide);
  \draw[arr] (slide) -- (pad);
  \draw[arr] (pad)   -- (split);
  \draw[arr] (split) -- (data);
  \draw[arr] (data)  -- (model);
  \draw[arr] (model) -- (train);
\end{tikzpicture}
\caption{End-to-end pipeline for the next-word prediction hands-on
  (\texttt{pytorch-lstm-next-word-predictor.ipynb}).}
\label{fig:pipeline}
\end{figure}

\subsection{Step 1: Tokenisation and Vocabulary Construction}
\label{subsec:tokenisation}

The raw text is processed as follows:

\begin{enumerate}[label=\textbf{\arabic*.}]
  \item \textbf{Lowercase conversion.}  All characters are converted to
    lowercase so that \texttt{About} and \texttt{about} are treated as the
    same token.  Python is case-sensitive; without this step the vocabulary
    would contain duplicates.

  \item \textbf{Word tokenisation.}  NLTK's \texttt{word\_tokenize} splits
    the text into individual word tokens and punctuation, automatically
    discarding HTML-style artefacts (e.g.\ \verb|\/|).

  \item \textbf{Frequency counting.}  A \texttt{Counter} object maps each
    unique token to its frequency.  The keys of this dictionary become the
    vocabulary.

  \item \textbf{Index assignment.}  Starting from index 1 (index 0 is
    reserved for the \texttt{<UNK>} token to handle out-of-vocabulary words
    at test time), each new token receives the next integer index.
\end{enumerate}

The resulting vocabulary size is $V = 289$ unique tokens.

\subsection{Step 2: Converting Sentences to Index Sequences}
\label{subsec:index-seq}

The document is split on newline characters (\verb|\n|), yielding 78
sentences of varying length.  A helper function \texttt{text\_to\_indices}
converts each sentence from a list of string tokens to a list of integer
indices by looking up the vocabulary dictionary; unknown tokens receive
index 0.

\subsection{Step 3: Self-Supervised Sliding-Window Sequences}
\label{subsec:sliding-window}

Next-word prediction is a \emph{self-supervised} task: the supervision signal
is derived from the data itself without manual annotation.  For each sentence
of length $L$ (represented as indices $[w_1, w_2, \ldots, w_L]$), the
following training pairs are generated:

\begin{center}
\begin{tabular}{lll}
\toprule
\textbf{Input sequence $X$}   & \textbf{Target $y$} & \textbf{Interpretation} \\
\midrule
$[w_1]$                       & $w_2$               & Given first word, predict second \\
$[w_1, w_2]$                  & $w_3$               & Given first two, predict third \\
$[w_1, w_2, w_3]$             & $w_4$               & \ldots \\
$\vdots$                      & $\vdots$            & \\
$[w_1, \ldots, w_{L-1}]$      & $w_L$               & Given all but last, predict last \\
\bottomrule
\end{tabular}
\end{center}

This produces a total of 942 training sequences from the 78 sentences.
Sequence lengths vary from 2 tokens (shortest) up to 62 tokens (longest).

\begin{example}[title=Sliding-Window Output for Sentence 10]
\begin{verbatim}
training_sequence[10]
# [ [3, 7],
#   [3, 7, 5],
#   [3, 7, 5, 2],
#   [3, 7, 5, 2, 6],
#   [3, 7, 5, 2, 6, 7],
#   ... ]
\end{verbatim}
Each sub-list is one training example.  The last element of the sub-list is
the target word; everything before it is the input context.
\end{example}

\subsection{Step 4: Padding to Uniform Length}
\label{subsec:padding}

Because PyTorch processes data in \emph{batches}, all sequences within a
batch must share the same length.  The maximum sequence length found in the
dataset is 62, so all shorter sequences are left-padded with zeros:

\[
  \text{padded\_seq} = [\underbrace{0,\ldots,0}_{62-L},\; w_1,\; w_2,\; \ldots,\; w_L].
\]

Left-padding (rather than right-padding) is used so that the target word
always occupies the \emph{last} position of the padded array, enabling a
uniform \texttt{y = sequence[-1]} split.

\begin{warningbox}[title=Why Left-Pad, Not Right-Pad]
  If zeros were appended to the right, the last column of the padded matrix
  would contain zeros for all short sequences, and those zeros would become
  the prediction targets.  Training on zero-targets would teach the model to
  predict ``unknown'' at the end of every sentence --- a catastrophic failure
  of data preparation.  Left-padding avoids this entirely.
\end{warningbox}

\subsection{Step 5: Input/Label Split}
\label{subsec:xy-split}

After padding, each training example is a length-62 integer vector.  The
split is:

\begin{itemize}
  \item $X = \text{sequence}[0:61]$ (first 61 tokens) --- the input context.
  \item $y = \text{sequence}[61]$ --- the target word to predict.
\end{itemize}

This yields a dataset $X \in \mathbb{Z}^{942 \times 61}$ and
$y \in \mathbb{Z}^{942}$.

\subsection{Step 6: PyTorch Dataset and DataLoader}
\label{subsec:dataset}

A standard PyTorch \texttt{Dataset} subclass wraps $X$ and $y$, and a
\texttt{DataLoader} creates mini-batches of size 32 with shuffling:

\begin{itemize}
  \item \texttt{\_\_len\_\_}: returns 942 (total training examples).
  \item \texttt{\_\_getitem\_\_(i)}: returns $(X[i], y[i])$ as tensors.
\end{itemize}

The DataLoader draws random batches of 32 examples, so each epoch consists
of $\lceil 942 / 32 \rceil = 30$ gradient updates.

\subsection{Step 7: LSTMModel Architecture}
\label{subsec:model}

The model is implemented as a PyTorch \texttt{nn.Module} with three components
arranged in sequence:

\begin{center}
\begin{tabular}{llll}
\toprule
\textbf{Layer}          & \textbf{Class}      & \textbf{Input shape}        & \textbf{Output shape} \\
\midrule
Embedding               & \texttt{nn.Embedding}  & $(B, 61)$ indices           & $(B, 61, 100)$ \\
LSTM                    & \texttt{nn.LSTM}       & $(B, 61, 100)$              & $(B, 61, 150)$ \\
Linear (FC head)        & \texttt{nn.Linear}     & $(B, 150)$ (last step only) & $(B, 289)$ logits \\
\bottomrule
\end{tabular}
\end{center}

\begin{example}[title=LSTMModel Class (PyTorch)]
\begin{verbatim}
class LSTMModel(nn.Module):
    def __init__(self, vocab_size, embedding_dim,
                 hidden_size, num_layers=1):
        super(LSTMModel, self).__init__()
        self.embedding = nn.Embedding(vocab_size, embedding_dim)
        self.lstm = nn.LSTM(embedding_dim, hidden_size,
                            num_layers=num_layers,
                            batch_first=True)
        self.fc = nn.Linear(hidden_size, vocab_size)
        self.total_hidden_states = num_layers * hidden_size

    def forward(self, x):
        embedded = self.embedding(x)          # (B, T, d_e)
        lstm_out, (hidden_state, cell_state) \
            = self.lstm(embedded)             # (B, T, H)
        output = self.fc(lstm_out[:, -1, :])  # (B, V)
        return output

model = LSTMModel(
    vocab_size=100,
    embedding_dim=100,
    hidden_size=100,
    num_layers=1
)
\end{verbatim}
The line \texttt{lstm\_out[:, -1, :]} extracts only the hidden state at the
\emph{last} time step, which encodes the full context of the input sequence.
This single 150-dimensional vector is then projected to logits over the
289-word vocabulary.
\end{example}

\begin{keyinsight}[title=Why \texttt{batch\_first=True}?]
  By default, PyTorch's \texttt{nn.LSTM} expects input of shape
  $(T, B, d_e)$ (time first).  Setting \texttt{batch\_first=True} changes
  this to $(B, T, d_e)$ (batch first), which is more natural when indexing
  with \texttt{lstm\_out[:, -1, :]} (select all batches, last time step).
  Without this flag, the equivalent slice would be
  \texttt{lstm\_out[-1, :, :]} (last time step, all batches), which is less
  readable.
\end{keyinsight}

\subsection{Step 8: Training Loop}
\label{subsec:training}

The model is trained for 10 epochs with:
\begin{itemize}
  \item \textbf{Loss function:} Cross-entropy loss
    $\mathcal{L} = -\sum_{k=1}^{V} y_k \log \hat{p}_k$, applied to the
    vocabulary logits.
  \item \textbf{Optimiser:} Adam with learning rate $\eta = 0.01$.
  \item \textbf{Device:} GPU if available (\texttt{cuda}), else CPU.
\end{itemize}

\newpage
%% ============================================================
\section{Comparison Tables}
\label{sec:tables}
%% ============================================================

\subsection{LSTM Gate Properties}
\label{subsec:gate-table}

\begin{table}[ht]
\centering
\caption{Summary of LSTM gate properties for a single time step.}
\label{tab:gate-props}
\begin{tabularx}{\textwidth}{lXXXX}
\toprule
\textbf{Gate}    & \textbf{Equation}
                 & \textbf{Activation}
                 & \textbf{Acts on}
                 & \textbf{Role} \\
\midrule
Forget $f_t$
  & $W_{xf}x_t + W_{hf}h_{t-1} + b_f$
  & Sigmoid $\in [0,1]$
  & $C_{t-1}$
  & Decides how much old memory to erase \\[4pt]
Input $i_t$
  & $W_{xi}x_t + W_{hi}h_{t-1} + b_i$
  & Sigmoid $\in [0,1]$
  & $\tilde{C}_t$
  & Decides how much new info to write \\[4pt]
Candidate $\tilde{C}_t$
  & $W_{xC}x_t + W_{hC}h_{t-1} + b_C$
  & Tanh $\in [-1,1]$
  & (new proposal)
  & Generates candidate update values \\[4pt]
Output $o_t$
  & $W_{xo}x_t + W_{ho}h_{t-1} + b_o$
  & Sigmoid $\in [0,1]$
  & $\tanh(C_t)$
  & Filters memory for short-term output \\
\bottomrule
\end{tabularx}
\end{table}

\subsection{Hyperparameters Used in the Hands-On}
\label{subsec:hyperparams}

\begin{table}[ht]
\centering
\caption{Hyperparameter values used in the hands-on notebook.}
\label{tab:hyperparams}
\begin{tabular}{lll}
\toprule
\textbf{Hyperparameter}    & \textbf{Value} & \textbf{Notes} \\
\midrule
Vocabulary size ($V$)      & 289            & Unique tokens after lowercasing \\
Embedding dimension ($d_e$)& 100            & \texttt{nn.Embedding} output size \\
Hidden size ($H$)          & 150            & LSTM neurons per gate \\
Number of LSTM layers      & 1              & Single-layer LSTM \\
Sequence length ($T$)      & 61             & After left-padding to 62, minus label \\
Batch size                 & 32             & Mini-batch gradient descent \\
Epochs                     & 10             & Full passes over training data \\
Learning rate ($\eta$)     & 0.01           & Adam optimiser \\
Max sequence length        & 62             & Determined from training data \\
Training sequences         & 942            & Generated by sliding-window \\
Sentences in corpus        & 78             & After splitting on newlines \\
\bottomrule
\end{tabular}
\end{table}

\subsection{Forget Gate Numerical Results}
\label{subsec:forget-table}

\begin{table}[ht]
\centering
\caption{Forget gate numerical walkthrough results (4-dimensional example).}
\label{tab:forget-nums}
\begin{tabular}{ccccc}
\toprule
\textbf{Dim.} & $z_f[k]$ & $f_t[k] = \sigma(z_f[k])$ & $C_{t-1}[k]$ & $f_t[k]\cdot C_{t-1}[k]$ \\
\midrule
0 & $1.2$  & $0.77$ & $0.9$ & $0.69$ \\
1 & $-0.5$ & $0.38$ & $0.2$ & $0.08$ \\
2 & $0.3$  & $0.57$ & $0.8$ & $0.46$ \\
3 & $2.0$  & $0.88$ & $0.4$ & $0.35$ \\
\bottomrule
\end{tabular}
\end{table}

%% ============================================================
\section{Summary and Conclusions}
\label{sec:summary}
%% ============================================================

Session 20 achieved two closely linked objectives:

\paragraph{Objective 1: Demystifying LSTM Mathematics.}
By working through a complete 4-dimensional numerical example, the session
gave students a concrete, calculable understanding of how the three LSTM
gates operate.  The key equations --- \cref{eq:full-ft} through
\cref{eq:full-ht} --- were derived step by step, with each intermediate vector
explicitly computed and interpreted.  The central insight is that LSTM resolves
the vanishing gradient problem by (a) maintaining a dedicated cell-state
highway $C_t$ updated only additively, and (b) using learnable, input-dependent
gates to modulate what is forgotten, what is added, and what is exposed.

\paragraph{Objective 2: Building a Working Language Model.}
The hands-on demonstrated that the same LSTM equations, when wrapped in a
\texttt{nn.Module} with an embedding layer and a linear projection head, can
be trained end-to-end on a next-word prediction task using standard PyTorch
utilities.  All stages of the NLP pre-processing pipeline were covered:
tokenisation, vocabulary construction, self-supervised sliding-window sequence
generation, left-padding for batched processing, and dataset/dataloader setup.

\begin{keyinsight}[title=Core Takeaways from Session 20]
\begin{enumerate}[label=\arabic*.]
  \item The LSTM cell state $C_t$ is a \textbf{long-term memory highway};
    the hidden state $h_t$ is a \textbf{context-filtered short-term output}.
  \item All internal vectors ($x_t$, $h_t$, $C_t$, $f_t$, $i_t$, $o_t$,
    $\tilde{C}_t$) must share the same dimension $H$.
  \item The forget gate multiplies $C_{t-1}$ by values in $[0,1]$;
    the input gate adds a fraction of the candidate $\tilde{C}_t$.
    Together, \cref{eq:full-cell} is the only update to long-term memory.
  \item In PyTorch, \texttt{nn.LSTM} encapsulates all gate computations.
    With \texttt{batch\_first=True}, input shape is
    $(B, T, d_e)$ and \texttt{lstm\_out[:, -1, :]} gives the final-step
    hidden state for classification or prediction.
  \item Batched processing of variable-length sequences requires
    \textbf{padding}.  Left-padding ensures the true target token always
    occupies the last position, avoiding zero-target training artefacts.
\end{enumerate}
\end{keyinsight}

\paragraph{Looking Ahead.}
Session 21 will introduce the Gated Recurrent Unit (GRU), which achieves
comparable performance to LSTM with approximately 25\% fewer parameters by
merging the cell state and hidden state into a single vector and replacing
the three separate gates with just two (reset and update).

\clearpage

\end{document}
